\documentclass[aps,preprint]{revtex4}%
\usepackage{amsfonts}
\usepackage{amsmath}
\usepackage{amssymb}
\usepackage{graphicx}%
\setcounter{MaxMatrixCols}{30}
%TCIDATA{OutputFilter=latex2.dll}
%TCIDATA{Version=5.50.0.2960}
%TCIDATA{CSTFile=revtex4.cst}
%TCIDATA{Created=Sunday, July 19, 2009 10:44:53}
%TCIDATA{LastRevised=Thursday, January 14, 2010 13:42:52}
%TCIDATA{<META NAME="GraphicsSave" CONTENT="32">}
%TCIDATA{<META NAME="SaveForMode" CONTENT="1">}
%TCIDATA{BibliographyScheme=Manual}
%TCIDATA{<META NAME="DocumentShell" CONTENT="Articles\SW\REVTeX 4">}
%BeginMSIPreambleData
\providecommand{\U}[1]{\protect\rule{.1in}{.1in}}
%EndMSIPreambleData
\newtheorem{theorem}{Theorem}
\newtheorem{acknowledgement}[theorem]{Acknowledgement}
\newtheorem{algorithm}[theorem]{Algorithm}
\newtheorem{axiom}[theorem]{Axiom}
\newtheorem{claim}[theorem]{Claim}
\newtheorem{conclusion}[theorem]{Conclusion}
\newtheorem{condition}[theorem]{Condition}
\newtheorem{conjecture}[theorem]{Conjecture}
\newtheorem{corollary}[theorem]{Corollary}
\newtheorem{criterion}[theorem]{Criterion}
\newtheorem{definition}[theorem]{Definition}
\newtheorem{example}[theorem]{Example}
\newtheorem{exercise}[theorem]{Exercise}
\newtheorem{lemma}[theorem]{Lemma}
\newtheorem{notation}[theorem]{Notation}
\newtheorem{problem}[theorem]{Problem}
\newtheorem{proposition}[theorem]{Proposition}
\newtheorem{remark}[theorem]{Remark}
\newtheorem{solution}[theorem]{Solution}
\newtheorem{summary}[theorem]{Summary}
\newenvironment{proof}[1][Proof]{\noindent\textbf{#1.} }{\ \rule{0.5em}{0.5em}}
\begin{document}
\preprint{HEP/123-qed}
\title[Short title for running header]{Long title that appears on the title page}
\author{First Author}
\affiliation{My Institution}
\author{Second Author}
\affiliation{My Institution}
\author{Third Author}
\affiliation{Other Institution}
\keywords{one two three}
\pacs{PACS number}

\begin{abstract}
Shell document for REV\TeX{} 4.

\end{abstract}
\volumeyear{year}
\volumenumber{number}
\issuenumber{number}
\eid{identifier}
\date[Date text]{date}
\received[Received text]{date}

\revised[Revised text]{date}

\accepted[Accepted text]{date}

\published[Published text]{date}

\startpage{101}
\endpage{102}
\maketitle
\tableofcontents


\section{Introduction}

The $N$-representability problem, first explicitly introduced by Husimi
\cite{Husimi} then Coulson \cite{Coulson} is based on the realization that
states of the $N$-electrons forming molecular and atomic systems are
determined by at most two electron interactions that in turn determine reduced
two particle states, represented by Second Order Reduced Density Operators
(SORDO's) that belong to the subset, $S_{N}^{2},$ of the set of postive trace
class operators acting in the Hilbert space $\mathcal{H}^{2},$ describing
$2$-electrons. These SORDO's can be represented by Second Order Reduced
Density Matrices (SORDM's) which are determined by
\begin{equation}
\eta_{2}=\frac{1}{2}\binom{r}{2}\left(  \binom{r}{2}+1\right)  \sim r^{4}%
\end{equation}
complex parameters, where $r$ is the size of the orthonormal basis set that
spans a finite dimensional approximation to $\mathcal{H}^{1},$ which should be
contrasted to the number $\eta_{N}$
\begin{equation}
\eta_{N}=\frac{1}{2}\binom{r}{N}\left(  \binom{r}{N}+1\right)  \sim r^{2N}%
\end{equation}
of complex parameters that is needed in general to determine the density
matrix of an $N$-electron states which are representations of postive
normalized trace class operators acting in the Hilbert space, $\mathcal{H}%
^{N}$. A crucial difference is that $\eta_{2}$ does not depend on the number
of electrons while $\eta_{N}$ does so
\begin{equation}
\frac{\eta_{N}}{\eta_{2}}\sim r^{2N-4},
\end{equation}
which increase rapidly as the number of electrons increases.

The formulation of the problem can be extended to distinguish between pure
states, ensemble (also known as mixed) states and states that describe varying
number of particles.

Even though the significant physical interactions that determine the nature of
the state are most pairwise the fact that electrons are fermions implies a
basic underlying exchange symmetry [thus a conservation principle and an
invariance group] that can be thought of as an $N$-electron interaction. This
basic symmetry interaction determines a class of possible electron states from
which a physical state is determined by coulombic interaction between
electrons and electrons and nuclei, the kinetic energy of the electrons and
the presence of electric and magnetic fields. States determined by the
$\eta_{2}$ parameters are called reduced states in contrast to those
determined by those determined by the $\eta_{N}$ parameters. Every state
determines a reduced state however the set of conditions that intrinsically
characterize reduced states are not known.

We formulate the problem of determing whether an element of the set of postive
trace class operators, $\mathcal{B}_{1+}\left(  \mathcal{H}^{2}\right)  ,$
acting in $\mathcal{H}^{2}$ actually is a reduced state by considering whether
the bounded postive linear functional that it determines on the set of bounded
operators, $\mathcal{B}\left(  \mathcal{H}^{2}\right)  ,$ acting in
$\mathcal{H}^{2}$ can be extended to a positive bounded linear functional
acting on the set of bounded operators acting in $\mathcal{H}^{N}$ i.e. we
formulate the conditions for a postive trace class operator acting in
$\mathcal{H}^{2}$ to be $N$-representable in terms of whether it can be
extended to a bounded positive linear functional on $\mathcal{B}\left(
\mathcal{H}^{N}\right)  .$

\section{Background}

\subsection{Fermi Fock Space and Algebra}

We consider the fermion Fock space, $\mathcal{H}_{F},$ based on the one
particle space Hilbert space $\mathcal{H}^{1}$ and the fermion Fock algebra
$\mathcal{F}$, defined by
\begin{align}
\mathcal{H}_{\mathcal{F}}  &  =\underset{0\leq N\leq\infty}{\oplus}%
\mathcal{H}^{N},\\
\mathcal{H}^{N}  &  =\wedge^{N}\mathcal{H}^{1}%
\end{align}
and
\begin{equation}
\mathcal{F}=\mathcal{B}\left(  \mathcal{H}_{\mathcal{F}}\right)  ,
\end{equation}
where $\wedge$ denotes antisymmetric tensor product and $\mathcal{B}\left(
\mathcal{H}_{\mathcal{F}}\right)  $ is the space of bounded operators acting
in $\mathcal{H}_{\mathcal{F}}$. The Fock algebra can be generated by
polynomials of second quantized operators, $\left\{  \left.  a_{j}^{\dagger
}\right\vert 1\leq j\leq\infty\right\}  ,$ which are defined by their action
on the vacuum state $\left\vert \phi\right\rangle $ by
\begin{equation}
a_{j}^{\dagger}\left\vert \phi\right\rangle =\left\vert \varphi_{j}%
\right\rangle ,1\leq j\leq\infty,
\end{equation}
where $\left\{  \left.  \left\vert \varphi_{j}\right\rangle \right\vert 1\leq
j\leq\infty\right\}  $ is a complete orthonormal basis of one-particle space
$\mathcal{H}^{1}$ and satisfy the anti-commutation relationships given by
\begin{equation}
\left[  a_{j}^{\dagger},a_{k}\right]  _{+}=a_{j}^{\dagger}a_{k}+a_{k}%
a_{j}^{\dagger}=\delta_{jk};1\leq j,k\leq\infty.
\end{equation}
The operators $\left\{  \left.  a_{j,}a_{j}^{\dagger}\right\vert 1\leq
j\leq\infty\right\}  $ can be considered as the values of the maps
\begin{equation}
\Phi,\Phi^{\dagger}:\mathcal{H}^{1}\rightarrow\mathcal{B}\left(
\mathcal{H}_{\mathcal{F}}\right)
\end{equation}
given by
\begin{align}
a_{j}  &  =\Phi\left(  \varphi_{j}\right)  ,\\
a_{j}^{\dagger}  &  =\Phi\left(  \varphi_{j}\right)  ^{\dagger}\equiv
\Phi^{\dagger}\left(  \varphi_{j}\right)  .
\end{align}
The point field operators $\left\{  \left.  \Phi\left(  \boldsymbol{r}\right)
\right\vert \boldsymbol{r\in}\mathbb{R}^{3}\right\}  $ are defined by
\cite{Thirring}
\begin{equation}
\Phi\left(  \boldsymbol{r}\right)  =\sum_{1\leq j\leq\infty}\overline
{\varphi_{j}\left(  \boldsymbol{r}\right)  }a_{j}.
\end{equation}
The unbounded operators $\Phi\left(  \boldsymbol{r}\right)  $ are densely
defined (but not closeable) on the Fock space, $\mathcal{H}_{\mathcal{F}},$
nor the subspaces $\mathcal{H}^{N}=\wedge^{N}\mathcal{H}^{1},$ and the domains
of their adjoints $\Phi\left(  \boldsymbol{r}\right)  ^{\dagger}$ are empty
i.e. do not contain any vectors belonging to $\mathcal{H}_{\mathcal{F}}$
\cite{Thirring}. However, matrix elements, $\left\langle \varphi\left\vert
\Phi\left(  \boldsymbol{r}\right)  ^{\dagger}\Phi\left(  \boldsymbol{r}%
^{\prime}\right)  \right.  \psi\right\rangle ,$ can be defined as the
distributional limits of $\left\langle \varphi\left\vert \Phi\left(  f\right)
^{\dagger}\Phi\left(  f^{\prime}\right)  \right.  \psi\right\rangle $ as
$f,f^{\prime}\rightarrow$ delta functions.

\subsection{Ket-Bra Representation of Second Quantized Operators}

\subsubsection{First Quantization of Second Quantization Operators}

\bigskip By considering the restrictions of $a_{j}^{\dagger}$ to
$\mathcal{H}^{N}$ for $0\leq N\leq\infty$ one obtains%
\begin{align}
a_{j}^{\dagger}  &  =\left\vert \varphi_{j}\right\rangle \left\langle
\phi\right\vert +\sum_{1\leq N\leq\infty}\sum_{1\leq j_{2}<\cdots<j_{N}%
\leq\infty}\left\vert \varphi_{j}\varphi_{j_{2}}\cdots\varphi_{j_{N}%
}\right\rangle \left\langle \varphi_{j_{2}}\cdots\varphi_{j_{N}}\right\vert
\nonumber\\
a_{j}  &  =\left\vert \phi\right\rangle \left\langle \varphi_{j}\right\vert
+\sum_{1\leq N\leq\infty}\sum_{1\leq j_{2}<\cdots<j_{N}\leq\infty}\left\vert
\varphi_{j_{2}}\cdots\varphi_{j_{N+1}}\right\rangle \left\langle \varphi
_{j}\varphi_{j_{2}}\cdots\varphi_{j_{N+1}}\right\vert . \label{FQ}%
\end{align}
These expansions \emph{do }preserve the algebraic structure as can be observed
from%
\begin{align}
a_{j}^{\dagger}a_{k}  &  =\sum_{\substack{1\leq N,M\leq\infty\\1\leq
j_{2}<\cdots<j_{N}\leq\infty\\1\leq k_{2}<\cdots<k_{M}\leq\infty}}\left\vert
\varphi_{j}\varphi_{j_{2}}\cdots\varphi_{j_{N}}\right\rangle \left\langle
\left.  \varphi_{j_{2}}\cdots\varphi_{j_{N}}\right\vert \varphi_{k_{2}}%
\cdots\varphi_{k_{M}}\right\rangle \left\langle \varphi_{k}\varphi_{k_{2}%
}\cdots\varphi_{k_{M}}\right\vert \nonumber\\
&  =\frac{1}{\left(  N-1\right)  !}\sum_{\sigma\in\mathcal{S}_{N-1}}\left(
-1\right)  ^{\pi\left(  \sigma\right)  }\delta_{j_{2}k_{\sigma\left(
2\right)  }}\cdots\delta_{j_{N}k_{\sigma\left(  N\right)  }}\sum
_{\substack{1\leq N\leq\infty\\1\leq j_{2}<\cdots<j_{N}\leq\infty\\1\leq
k_{2}<\cdots<k_{N}\leq\infty}}\left\vert \varphi_{j}\varphi_{j_{2}}%
\cdots\varphi_{j_{N}}\right\rangle \left\langle \varphi_{k}\varphi_{k_{2}%
}\cdots\varphi_{k_{N}}\right\vert ,
\end{align}


which leads to
\begin{equation}
a_{j}^{\dagger}a_{k}=\sum_{1\leq N\leq\infty}\sum_{1\leq j_{2}<\cdots
<j_{N}\leq\infty}\left\vert \varphi_{j}\varphi_{j_{2}}\cdots\varphi_{j_{N}%
}\right\rangle \left\langle \varphi_{k}\varphi_{j_{2}}\cdots\varphi_{j_{N}%
}\right\vert ,
\end{equation}
\begin{equation}
a_{j}^{\dagger}a_{j}=\sum_{1\leq N\leq\infty}\sum_{1\leq j_{2}<\cdots
<j_{N}\leq\infty}\left\vert \varphi_{j}\varphi_{j_{2}}\cdots\varphi_{j_{N}%
}\right\rangle \left\langle \varphi_{j}\varphi_{j_{2}}\cdots\varphi_{j_{N}%
}\right\vert ,
\end{equation}
and
\begin{equation}
\mathcal{N=}\sum_{1\leq j\leq\infty}a_{j}^{\dagger}a_{j}=\sum_{1\leq
N\leq\infty}NP_{\mathcal{H}^{N}}. \label{NumOp}%
\end{equation}
Eq. $\left(  \ref{NumOp}\right)  $ can be seen by considering the example%
\begin{equation}
\sum_{1\leq j\leq\infty}\left.  a_{j}^{\dagger}a_{j}\right\vert _{\mathcal{H}%
^{2}}=\sum_{1\leq j,j_{2}\leq\infty}\left\vert \varphi_{j}\varphi_{j_{2}%
}\right\rangle \left\langle \varphi_{j}\varphi_{j_{2}}\right\vert
=2\sum_{1\leq j<j_{2}\leq\infty}\left\vert \varphi_{j}\varphi_{j_{2}%
}\right\rangle \left\langle \varphi_{j}\varphi_{j_{2}}\right\vert
=2P_{\mathcal{H}^{2}}.
\end{equation}


\subsubsection{Second Quantization of First Quantized Operators}

The inverse relationship between FQ and SQ operators is%

\begin{equation}
\left\vert \varphi_{j_{1}}\cdots\varphi_{j_{N}}\right\rangle \left\langle
\varphi_{k_{1}}\cdots\varphi_{k_{N}}\right\vert =P_{N}a_{j_{1}}^{\dagger
}\cdots a_{j_{N}}^{\dagger}a_{k_{N}}\cdots a_{k_{1}}P_{N}%
\end{equation}


\subsection{Particle Conserving Operators}

The operator ideal spaces $\left\{  \mathcal{C}\left(  \mathcal{H}%
_{\mathcal{F}}\right)  ,\mathcal{B}_{k}\left(  \mathcal{H}_{\mathcal{F}%
}\right)  ;1\leq\kappa\leq2\right\}  $ are contained in $\mathcal{B}\left(
\mathcal{H}_{\mathcal{F}}\right)  $ and
\begin{equation}
\mathcal{B}\left(  \mathcal{H}_{\mathcal{F}}\right)  \supseteq\mathcal{C}%
\left(  \mathcal{H}_{\mathcal{F}}\right)  \supseteq\mathcal{B}_{2}\left(
\mathcal{H}_{\mathcal{F}}\right)  \supseteq\mathcal{B}_{1}\left(
\mathcal{H}_{\mathcal{F}}\right)  \supseteq\mathcal{B}_{0}\left(
\mathcal{H}_{\mathcal{F}}\right)  ,
\end{equation}
where $\mathcal{C}\left(  \mathcal{H}_{\mathcal{F}}\right)  $ are compact,
$\mathcal{B}_{2}\left(  \mathcal{H}_{\mathcal{F}}\right)  $ are Hilbert
Schmidt, $\mathcal{B}_{1}\left(  \mathcal{H}_{\mathcal{F}}\right)  $ are trace
class and $\mathcal{B}_{0}\left(  \mathcal{H}_{\mathcal{F}}\right)  $ are
finite rank operators. The spaces $\left\{  \mathcal{C}\left(  \mathcal{H}%
_{\mathcal{F}}\right)  ,\mathcal{B}_{k}\left(  \mathcal{H}_{\mathcal{F}%
}\right)  ;1\leq\kappa\leq2\right\}  $ are normed spaces, which are the
completions of $\mathcal{B}_{0}\left(  \mathcal{H}_{\mathcal{F}}\right)  $ in
the appropiate norms. Thus it is often more convenient to work within the
context of finite rank operators and take the appropiate limits (if they
exist) when necessary.

The self adjoint number operator, $\mathcal{N},$ defined as
\begin{equation}
\mathcal{N=}\sum_{1\leq j\leq\infty}a_{j}^{\dagger}a_{j}=\int\Phi\left(
\boldsymbol{x}\right)  ^{\dagger}\Phi\left(  \boldsymbol{x}\right)  dx
\end{equation}
allows one to decomposes the space $\mathcal{B}\left(  \mathcal{H}%
_{\mathcal{F}}\right)  $ into a direct sum of operators that commute with
$\mathcal{N}$ and those that do not as
\begin{equation}
\mathcal{B}\left(  \mathcal{H}_{\mathcal{F}}\right)  =\mathcal{B}\left(
\mathcal{H}_{\mathcal{F}}\right)  _{e}\oplus\mathcal{B}\left(  \mathcal{H}%
_{\mathcal{F}}\right)  _{o}.
\end{equation}
The subspace (which is also a subalgebra) $\mathcal{B}\left(  \mathcal{H}%
_{\mathcal{F}}\right)  _{e}$ can be identified with polynomials and strong
limits of polynomials of even numbers of creators and annihilators in
\emph{normal form} i.e.%
\begin{equation}
X=\sum_{\substack{1\leq P\leq\infty\\1\leq j_{1}<\cdots<j_{P}\leq\infty\\1\leq
k_{1}<\cdots<k_{P}\leq\infty}}X_{j_{1}\cdots j_{P}k_{1}\cdots k_{P}}a_{j_{1}%
}^{\dagger}\cdots a_{j_{P}}^{\dagger}a_{k_{P}}\cdots a_{k_{1}},
\end{equation}
while the subspace $\mathcal{B}\left(  \mathcal{H}_{\mathcal{F}}\right)  _{o}$
corresponds to an odd number. The subspace $\mathcal{B}\left(  \mathcal{H}%
_{\mathcal{F}}\right)  _{e}$ decomposes as the direct sums
\begin{align}
\mathcal{F}_{e} &  =\mathcal{B}\left(  \mathcal{H}_{\mathcal{F}}\right)
_{e}\nonumber\\
&  =%
%TCIMACRO{\dbigoplus \limits_{0\leq N\leq\infty}}%
%BeginExpansion
{\displaystyle\bigoplus\limits_{0\leq N\leq\infty}}
%EndExpansion
\mathcal{B}\left(  \mathcal{H}_{\mathcal{F}}\right)  _{eN}\label{BanSp}\\
&  \cong%
%TCIMACRO{\dbigoplus \limits_{0\leq N\leq\infty}}%
%BeginExpansion
{\displaystyle\bigoplus\limits_{0\leq N\leq\infty}}
%EndExpansion
\mathcal{B}\left(  \mathcal{H}^{N}\right)  ,\label{BanAlg}%
\end{align}
where the two Banach space direct sums are \emph{different} but isomorphic. In
order to show the validity of the first decomposition one needs to show that
\begin{equation}
\mathcal{B}\left(  \mathcal{H}_{\mathcal{F}}\right)  _{eN}\cap\mathcal{B}%
\left(  \mathcal{H}_{\mathcal{F}}\right)  _{eN^{\prime}}=\phi
\end{equation}
or equivalently that
\begin{equation}
\sum_{0\leq N\leq\infty}A_{N}=0\Leftrightarrow A_{N}=0,0\leq N\leq\infty
,A_{N}\in\mathcal{B}_{sa}\left(  \mathcal{H}^{N}\right)  .
\end{equation}
It is easy to see that the last condition is satisfied, consider
\begin{equation}
\left(  \sum_{0\leq N\leq\infty}A_{N}\right)  \left\vert \Psi_{M}%
\right\rangle
\end{equation}
and chose $\left\vert \Psi_{M}\right\rangle \in\mathcal{H}^{M},$ $0\leq
M\leq\infty.$

The spaces $\mathcal{B}_{0sa}\left(  \mathcal{H}^{N}\right)  ,0\leq
N\leq\infty$ are othogonal with respect to the H-S norm, but $\mathcal{B}%
_{0sa}\left(  \mathcal{H}_{\mathcal{F}}\right)  _{eN}$ are not, consider%
\begin{align}
&  \sum_{\substack{1\leq j_{1}<\cdots<j_{P}\leq\infty\\1\leq k_{1}%
<\cdots<k_{P}\leq\infty}}X_{j_{1}\cdots j_{P}k_{1}\cdots k_{P}}%
\operatorname*{Tr}\left\{  a_{j_{1}}^{\dagger}\cdots a_{j_{P}}^{\dagger
}a_{k_{P}}\cdots a_{k_{1}}\mathcal{N}_{1}\right\}  ,P\neq N\nonumber\\
&  =\sum_{\substack{1\leq j_{1}<\cdots<j_{P}\leq\infty\\1\leq k_{1}%
<\cdots<k_{P}\leq\infty}}X_{j_{1}\cdots j_{P}k_{1}\cdots k_{P}}%
P\operatorname*{Tr}\nolimits_{N}\left\{  a_{j_{1}}^{\dagger}\cdots a_{j_{P}%
}^{\dagger}a_{k_{P}}\cdots a_{k_{1}}\right\}  \neq0.
\end{align}


The subspaces $\left\{  \left.  \mathcal{B}\left(  \mathcal{H}^{N}\right)
\right\vert 0\leq N\leq\infty\right\}  $ are subalgebras while the subspaces
$\left\{  \left.  \mathcal{B}\left(  \mathcal{H}_{\mathcal{F}}\right)
_{eN}\right\vert 0\leq N\leq\infty\right\}  $ \emph{are not }(the identity has
components in all the spaces and $\mathcal{B}\left(  \mathcal{H}_{\mathcal{F}%
}\right)  _{e1}$ has to have the identity added to it to be closed). Thus eq.
$\left(  \ref{BanAlg}\right)  $ is a direct sum of algebras while eq. $\left(
\ref{BanSp}\right)  $ is just a direct sum of banach spaces.

\subsection{Complexification}

As density operators and hamiltonians are in general self adjoint we can
restrict attention to the real space $\mathcal{B}_{sa}\left(  \mathcal{H}%
_{\mathcal{F}}\right)  ,$which is a real subspace of $\mathcal{B}\left(
\mathcal{H}_{\mathcal{F}}\right)  $ such thats its complexification gives
$\mathcal{B}\left(  \mathcal{H}_{\mathcal{F}}\right)  $ i.e.
\begin{equation}
\mathcal{B}\left(  \mathcal{H}_{\mathcal{F}}\right)  =\mathcal{B}_{sa}\left(
\mathcal{H}_{\mathcal{F}}\right)  \oplus i\mathcal{B}_{sa}\left(
\mathcal{H}_{\mathcal{F}}\right)
\end{equation}
such that $Z\in\mathcal{B}_{2}\left(  \mathcal{H}^{N}\right)  ,$ where
\begin{equation}
Z=X+iY;\text{ }X,Y\in\mathcal{B}_{2sa}\left(  \mathcal{H}^{N}\right)  .
\end{equation}
The inner product in $\mathcal{B}_{2}\left(  \mathcal{H}^{N}\right)  $ is
defined in terms of the inner product in $\mathcal{B}_{2sa}\left(
\mathcal{H}^{N}\right)  $ as
\begin{equation}
\left(  Z_{1}\left\vert Z_{2}\right.  \right)  _{HS}=\left(  X_{1}\left\vert
X_{2}\right.  \right)  +\left(  Y_{1}\left\vert Y_{2}\right.  \right)
+i\left\{  \left(  X_{1}\left\vert Y_{2}\right.  \right)  -\left(
Y_{1}\left\vert X_{2}\right.  \right)  \right\}  ,
\end{equation}
which corresponds to
\begin{equation}
\left(  Z_{1}\left\vert Z_{2}\right.  \right)  _{HS}=\operatorname{Tr}\left\{
Z_{1}^{\dagger}Z_{2}\right\}  ,
\end{equation}
as%
\begin{equation}
\operatorname{Tr}\left\{  \left(  X_{1}-iY_{1}\right)  \left(  X_{2}%
+iY_{2}\right)  \right\}  =\operatorname{Tr}\left\{  X_{1}X_{2}-iY_{1}%
X_{2}+iX_{1}Y_{2}+Y_{1}Y_{2}\right\}  .
\end{equation}


For density operators and hamiltonians that commute with the number operator
$\mathcal{N}$ we can wlog consider $\mathcal{B}_{sa}\left(  \mathcal{H}%
_{\mathcal{F}}\right)  _{e}.$

Note that $\mathcal{B}_{2sa}\left(  \mathcal{H}_{\mathcal{F}}\right)  $ (self
adjoint operators) and $i\mathcal{B}_{2sa}\left(  \mathcal{H}_{\mathcal{F}%
}\right)  $ (anti self adjoint operators) are not orthogonal subspaces of
$\mathcal{B}_{2}\left(  \mathcal{H}_{\mathcal{F}}\right)  $ as they are not
even subspaces under $\mathbb{C}$-linear combination.

The are not closed under the operator product, however they are closed under
the Jordan product and thus form Jordan algebras (which are non associative).

\section{Expansion and Contraction Maps.}

The contraction maps, $\left\{  \left.  L_{N}^{P}\right\vert 0\leq P\leq
N;0\leq N\leq\infty\right\}  ,$ map $N$-particle states%
\begin{equation}
S_{N}=\left\{  D^{N}\left\vert D^{N}\in\mathcal{B}_{1+}\left(  \mathcal{H}%
^{N}\right)  ;\operatorname{Tr}D^{N}=1\right.  \right\}  ,
\end{equation}
to reduced $P$-particle states
\begin{equation}
L_{N}^{P}:S_{N}\subset\mathcal{B}_{1+}\left(  \mathcal{H}^{N}\right)
\subset\mathcal{B}_{1sa}\left(  \mathcal{H}^{N}\right)  \longrightarrow
\tilde{S}_{NP}\subset S_{P}\subset\mathcal{B}_{1+}\left(  \mathcal{H}%
^{P}\right)  \subset\mathcal{B}_{1sa}\left(  \mathcal{H}^{P}\right)  ,
\end{equation}
where $\mathcal{B}_{1sa}\left(  \mathcal{H}^{N}\right)  $ is the space of self
adjoint trace class operators acting in $\mathcal{H}^{N},$ $\mathcal{B}%
_{1+}\left(  \mathcal{H}^{N}\right)  $ the cone of positive self adjoint trace
class operators, $S_{N}$ the convex set of $N$-particle states, $S_{P}%
,\mathcal{B}_{1+}\left(  \mathcal{H}^{P}\right)  ,\mathcal{B}_{1sa}\left(
\mathcal{H}^{P}\right)  $ the corresponding $P$-particle objects and
$\tilde{S}_{NP}$ the set of $N$- representable $P^{th}$ order Reduced Density
Operators (PRDO's) that are normalized to $1$ i.e.
\begin{equation}
\tilde{S}_{NP}=\left\{  \left.  \tilde{D}_{NP}\right\vert \tilde{D}_{NP}%
=L_{N}^{P}\left(  D^{N}\right)  ,D^{N}\in S_{N}\right\}  .
\end{equation}
The set, $\left\{  D_{NP}\right\}  ,$ of reduced density operators
\begin{equation}
D_{NP}=\binom{N}{P}\tilde{D}_{NP},
\end{equation}
normalized to the number of $P$-particle groups is denoted by $S_{NP}$ i.e.
\begin{equation}
S_{NP}=\left\{  D_{NP}\left\vert D_{NP}=\binom{N}{P}L_{N}^{P}\left(
D^{N}\right)  ,D^{N}\in S_{N}\right.  \right\}  .
\end{equation}
The contraction maps are explicitly defined by
\begin{align}
\tilde{D}_{NP} &  =L_{N}^{P}\left(  D^{N}\right)  =\binom{N}{P}^{-1}%
\sum_{\substack{1\leq j_{1}<\cdots<j_{P}\leq\infty\\1\leq k_{1}<\cdots
<k_{P}\leq\infty}}\operatorname{Tr}\left\{  a_{j_{1}}^{\dagger}\cdots
a_{j_{P}}^{\dagger}a_{k_{P}}\cdots a_{k_{1}}D^{N}\right\}  \left\vert
\varphi_{j_{1}}\cdots\varphi_{j_{P}}\right\rangle \left\langle \varphi_{k_{1}%
}\cdots\varphi_{k_{P}}\right\vert \nonumber\\
&  =\binom{N}{P}^{-1}\sum_{\substack{1\leq j_{1}<\cdots<j_{P}\leq\infty\\1\leq
k_{1}<\cdots<k_{P}\leq\infty}}D_{NPj_{1}\cdots j_{P}k_{1}\cdots k_{P}%
}\left\vert \varphi_{j_{1}}\cdots\varphi_{j_{P}}\right\rangle \left\langle
\varphi_{k_{1}}\cdots\varphi_{k_{P}}\right\vert ;0\leq P\leq N.
\end{align}


As $S_{N},S_{P},S_{NP}$ and $\tilde{S}_{NP}$ are convex sets they are
generated by their extreme points $\operatorname*{ext}\left\{  S_{N}\right\}
,\operatorname*{ext}\left\{  S_{P}\right\}  ,\operatorname*{ext}\left\{
S_{NP}\right\}  $ and $\operatorname*{ext}\left\{  \tilde{S}_{NP}\right\}
.$The extreme points $\operatorname*{ext}\left\{  S_{N}\right\}
,\operatorname*{ext}\left\{  S_{P}\right\}  $ of $S_{N}$ and $S_{P}$
respectively are pure states that are projectors onto one dimensional
subspaces of $\mathcal{H}^{N}$ and $\mathcal{H}^{P}$ (their rank is 1), which
should be constrasted to $\operatorname*{ext}\left\{  \tilde{S}_{NP}\right\}
$ whose rank can never be one ($N\neq P,0$ ) and thus must generate a subset
of $S_{P}.$These subsets are convex sets of \emph{mixed states}$.$ In the case
$P=1$
\begin{equation}
\tilde{S}_{N1}=\left\{  \left.  \tilde{D}_{N1}\right\vert \tilde{D}_{N1}%
=L_{N}^{1}\left(  D^{N}\right)  ,D^{N}\in S_{N}\right\} \nonumber
\end{equation}
It is easy to see that
\begin{equation}
\operatorname*{Tr}\left\{  \tilde{D}_{N1}\right\}  =1\text{ and }\tilde
{D}_{N1}\geq0\text{ for }\tilde{D}_{N1}\in\tilde{S}_{N1}%
\end{equation}
thus
\begin{equation}
\tilde{S}_{N1}\subset S_{1}.
\end{equation}
Convex sets are generated by their extreme points and one can show that
\begin{equation}
\operatorname*{ext}\left\{  \tilde{S}_{N1}\right\}  =\left\{  \tilde{D}%
_{N1}\left\vert \tilde{D}_{N1}=\frac{1}{N}\sum_{1\leq j\leq N}\left\vert
\psi_{j}\right\rangle \left\langle \psi_{j}\right\vert ;\text{ where }\left\{
\left.  \left\vert \psi_{j}\right\rangle \right\vert 1\leq j\leq
\infty\right\}  \text{ are conb's of }\mathcal{H}^{1}\right.  \right\}  ,
\end{equation}
thus%
\begin{equation}
\tilde{S}_{N1}=\left\{  \tilde{D}_{1}\left\vert \tilde{D}_{1}=\sum_{1\leq
j\leq\infty}\beta_{j}\tilde{D}_{N1j};\sum_{1\leq j\leq\infty}\beta_{j}%
=1,0\leq\beta_{j}\leq1,1\leq j\leq\infty\right.  \right\}  ,
\end{equation}
where $\tilde{D}_{N1j}\in\operatorname*{ext}\left\{  \tilde{S}_{N1}\right\}
,1\leq j\leq\infty.$ As $\tilde{S}_{N1}\subset S_{1}$ one has the spectral
decomposition
\begin{equation}
\tilde{S}_{N1}=\left\{  \tilde{D}_{1}\left\vert \tilde{D}_{1}=\sum_{1\leq
j\leq\infty}\alpha_{j}\left\vert \psi_{j}\right\rangle \left\langle \psi
_{j}\right\vert ;\sum_{1\leq j\leq\infty}\alpha_{j}=1,0\leq\alpha_{j}%
\leq1,1\leq j\leq\infty\right.  \right\}  .
\end{equation}


In the case $P=2$
\begin{align}
&  \tilde{S}_{N2}=\left\{  \left.  \tilde{D}_{N2}\right\vert \tilde{D}%
_{N2}=L_{N}^{2}\left(  D^{N}\right)  ,D^{N}\in S_{N}\right\} \nonumber\\
&  =\left\{  \tilde{D}_{2N}\left\vert \tilde{D}_{2N}=\sum_{1\leq j\leq\infty
}\alpha_{j}\left\vert \psi_{2j}\right\rangle \left\langle \psi_{2j}\right\vert
;\sum_{1\leq j\leq\infty}\alpha_{j}=1,0\leq\alpha_{j}\leq1,\right.  \right.
\nonumber\\
&  \left.  1\leq j\leq\infty,\left\{  \left.  \left\vert \psi_{2j}%
\right\rangle \right\vert 1\leq j\leq\infty\right\}  \text{ are conb's of
}\mathcal{H}^{2}\right\} \nonumber\\
&  =\left\{  \tilde{D}\left\vert \tilde{D}=\sum_{1\leq j\leq\infty}\beta
_{j}\tilde{D}_{j};\sum_{1\leq j\leq\infty}\beta_{j}=1,0\leq\beta_{j}%
\leq1,1\leq j\leq\infty\right.  \right\}
\end{align}
and
\begin{equation}
\tilde{S}_{N2}\subset S_{2}.
\end{equation}
If one knew the extreme points of $\tilde{S}_{N2}$ the problem would be solved.

The preceding definitions produce the normalization
\begin{align}
&  \operatorname{Tr}\left\{  D_{N1}\right\}  =\sum_{1\leq j_{1}\leq\infty
}\operatorname{Tr}\left\{  a_{j_{1}}^{\dagger}a_{j_{1}}D^{N}\right\}  \text{
}\left\langle \left.  \varphi_{j_{1}}\right\vert \varphi_{j_{1}}\right\rangle
_{1}=\operatorname{Tr}\left\{  \mathcal{N}D^{N}\right\}  =N\nonumber\\
&  \operatorname{Tr}\left\{  D_{N2}\right\}  =\sum_{1\leq j_{1}<j_{2}%
\leq\infty}\operatorname{Tr}\left\{  a_{j_{1}}^{\dagger}a_{j_{2}}^{\dagger
}a_{j_{2}}a_{j_{1}}D^{N}\right\}  \left\langle \left.  \varphi_{j_{1}}%
\varphi_{j_{2}}\right\vert \varphi_{j_{1}}\varphi_{j_{2}}\right\rangle
_{2}=\operatorname{Tr}\left\{  \mathcal{N}_{2}D^{N}\right\}  =\binom{N}%
{2}\nonumber\\
&  \qquad\qquad\qquad\quad\qquad\qquad\qquad\quad\vdots\nonumber\\
&  \operatorname{Tr}\left\{  D_{NP}\right\}  =\sum_{1\leq j_{1}<\cdots
<j_{P}\leq\infty}\operatorname{Tr}\left\{  a_{j_{1}}^{\dagger}\cdots a_{j_{P}%
}^{\dagger}a_{j_{P}}\cdots a_{j_{1}}D^{N}\right\}  \left\langle \left.
\varphi_{j_{1}}\cdots\varphi_{j_{P}}\right\vert \varphi_{j_{1}}\cdots
\varphi_{j_{P}}\right\rangle _{P}\nonumber\\
&  \qquad\qquad\qquad\qquad\qquad\qquad\qquad\qquad\qquad\qquad\qquad
\qquad\quad\left.  =\right.  \operatorname{Tr}\left\{  \mathcal{N}_{P}%
D^{N}\right\}  =\binom{N}{P}\nonumber\\
&  \qquad\qquad\qquad\quad\qquad\qquad\qquad\quad\vdots
\end{align}
and
\begin{equation}
\operatorname{Tr}\left\{  \tilde{D}_{NP}\right\}  =1\,\forall N,P.
\end{equation}


In the coordinate representation
\begin{equation}
D^{p}\left(  \boldsymbol{\kappa}_{p};\boldsymbol{\kappa}_{p}^{\prime}\right)
=\operatorname{Tr}\left\{  \Phi\left(  \boldsymbol{\kappa}_{p}^{\prime
}\right)  ^{\dagger}\Phi\left(  \boldsymbol{\kappa}_{p}\right)  D\right\}
=\int D^{p}\left(  \boldsymbol{\kappa}_{p};\boldsymbol{\kappa}_{p}^{\prime
}\right)  \left\vert \boldsymbol{\kappa}_{p}\right\rangle \left\langle
\boldsymbol{\kappa}_{p}^{\prime}\right\vert d\boldsymbol{\kappa}%
_{p}d\boldsymbol{\kappa}_{p}^{\prime}, \label{DenMatDef}%
\end{equation}
where
\begin{align}
\boldsymbol{\kappa}_{p}  &  \equiv\boldsymbol{x}_{1},\cdots,\boldsymbol{x}%
_{p},\\
\boldsymbol{x}_{j}  &  \equiv\boldsymbol{r}_{j}\boldsymbol{,\xi}_{j};1\leq
j\leq p.\nonumber
\end{align}
The operators $\Phi\left(  \boldsymbol{\kappa}_{p}\right)  ^{\dagger}$ are
defined as the product of the one-electrons operators $\Phi\left(
\boldsymbol{\kappa}_{1}\right)  ^{\dagger}$ i.e.%
\begin{equation}
\Phi\left(  \boldsymbol{\kappa}_{p}\right)  ^{\dagger}=\Phi\left(
\boldsymbol{x}_{1}\right)  ^{\dagger}\cdots\Phi\left(  \boldsymbol{x}%
_{p}\right)  ^{\dagger}.
\end{equation}


The adjoint maps, $\Gamma_{P}^{N},$ wrt to the scalar products (which are
generaliztions of the Hilbert-Schmidt inner products ---the spaces
$\mathcal{B}_{sa}\left(  \mathcal{H}^{N}\right)  \supseteq\mathcal{B}%
_{2sa}\left(  \mathcal{H}^{N}\right)  \supseteq\mathcal{B}_{1sa}\left(
\mathcal{H}^{N}\right)  $ could be thought of as a Gelfand Triple.)
\begin{align}
\text{ }\left\langle \left.  \bullet\right\vert \bullet\right\rangle _{N}  &
:\mathcal{B}_{sa}\left(  \mathcal{H}^{N}\right)  \times\mathcal{B}%
_{1sa}\left(  \mathcal{H}^{N}\right)  \rightarrow\mathbb{R}\\
\text{ }\left\langle \left.  \bullet\right\vert \bullet\right\rangle _{p}  &
:\mathcal{B}_{sa}\left(  \mathcal{H}^{p}\right)  \times\mathcal{B}%
_{1sa}\left(  \mathcal{H}^{p}\right)  \rightarrow\mathbb{R}%
\end{align}
are defined as
\begin{equation}
\left\langle \left.  \Gamma_{P}^{N}\left(  X\right)  \right\vert
Y\right\rangle _{N}=\left\langle \left.  X\right\vert L_{N}^{P}\left(
Y\right)  \right\rangle _{P}.
\end{equation}
The maps $\left\{  L_{N}^{P},\Gamma_{P}^{N}\right\}  $ can be naturally
extended to $\mathcal{B}\left(  \mathcal{H}_{\mathcal{F}}\right)
_{sae},\mathcal{B}\left(  \mathcal{H}_{\mathcal{F}}\right)  _{sa}$ and
$\mathcal{B}\left(  \mathcal{H}_{\mathcal{F}}\right)  .$ $L_{N}^{P}$ is
surjective, while $\Gamma_{P}^{N}$ is injective, thus in the infinite
dimensional case $\Gamma_{P}^{N}\left(  \mathcal{B}_{sa}\left(  \mathcal{H}%
^{P}\right)  \right)  $ is a proper subspace of $\mathcal{B}_{sa}\left(
\mathcal{H}^{N}\right)  $ for $N>P,$while in the finite dimensional case for
$P<\left[  \frac{N}{2}\right]  .$ Further as
\begin{equation}
\left\langle \left.  \Gamma_{P}^{N}\left(  X\right)  \right\vert
Y\right\rangle _{N}=\left\langle \left.  X\right\vert L_{N}^{P}\left(
Y\right)  \right\rangle _{P}=0,\,Y\in\ker\left\{  L_{N}^{P}\right\}
\end{equation}
one has
\begin{equation}
\operatorname*{ran}\left\{  \Gamma_{P}^{N}\right\}  \subset\mathcal{B}\left(
\mathcal{H}^{N}\right)  \bot\ker\left\{  L_{N}^{P}\right\}  \subset
\mathcal{B}_{1}\left(  \mathcal{H}^{N}\right)  .
\end{equation}
Taking the closure in Hilbert-Schmidt norm one obtains
\begin{equation}
\mathcal{B}_{2sa}\left(  \mathcal{H}^{N}\right)  =\overline{\ker\left\{
L_{N}^{P}\right\}  }\oplus\ker\left\{  L_{N}^{P}\right\}  ^{\bot}.
\end{equation}%
\begin{equation}
\ker\left\{  L_{N}^{P}\right\}  ^{\bot}=%
%TCIMACRO{\dbigcup \limits_{P\leq P^{\prime}<N}}%
%BeginExpansion
{\displaystyle\bigcup\limits_{P\leq P^{\prime}<N}}
%EndExpansion
\operatorname*{ran}\left\{  \Gamma_{P^{\prime}}^{N}\right\}
\end{equation}


We can define operators, $T_{NP},$ belonging to $\mathcal{B}_{1sa}\left(
\mathcal{H}_{\mathcal{F}}\right)  _{e}$ that act in $\mathcal{H}_{\mathcal{F}%
}$
\begin{equation}
\binom{N}{P}T_{NP}=\sum_{\substack{1\leq j_{1}<\cdots<j_{P}\leq\infty\\1\leq
k_{1}<\cdots<k_{P}\leq\infty}}\operatorname{Tr}\left\{  a_{j_{1}}^{\dagger
}\cdots a_{j_{P}}^{\dagger}a_{k_{P}}\cdots a_{k_{1}}D^{N}\right\}  a_{j_{1}%
}^{\dagger}\cdots a_{j_{P}}^{\dagger}a_{k_{P}}\cdots a_{k_{1}}%
\end{equation}
and in particular both $\mathcal{H}^{P^{\prime}}$ and $\mathcal{H}^{N}$ by
taking restrictions$.$ One can observe that
\begin{equation}
D_{NP^{\prime}}=L_{N}^{P^{\prime}}\left(  D^{N}\right)  =\left.
T_{NP}\right\vert _{\mathcal{H}^{P^{\prime}}};0\leq P^{\prime}\leq P;0\leq
P\leq N
\end{equation}
and prove that
\begin{equation}
\Gamma_{P}^{N}\left(  D_{NP}\right)  =\left.  T_{NP}\right\vert _{\mathcal{H}%
^{N}}.
\end{equation}


$\lceil$%
\begin{align}
T_{NN}  &  =\sum_{\substack{1\leq j_{1}<\cdots<j_{N}\leq\infty\\1\leq
k_{1}<\cdots<k_{N}\leq\infty}}\operatorname{Tr}\left\{  a_{j_{1}}^{\dagger
}\cdots a_{j_{N}}^{\dagger}a_{k_{N}}\cdots a_{k_{1}}D^{N}\right\}  a_{j_{1}%
}^{\dagger}\cdots a_{j_{N}}^{\dagger}a_{k_{N}}\cdots a_{k_{1}}\nonumber\\
\left.  T_{NN}\right\vert _{\mathcal{H}^{N}}  &  =D^{N}\nonumber\\
D_{NP}  &  =\binom{N}{P}L_{N}^{P}\left(  D^{N}\right)  =\sum_{\substack{1\leq
j_{1}<\cdots<j_{P}\leq\infty\\1\leq k_{1}<\cdots<k_{P}\leq\infty
}}\operatorname{Tr}\left\{  a_{j_{1}}^{\dagger}\cdots a_{j_{P}}^{\dagger
}a_{k_{P}}\cdots a_{k_{1}}D^{N}\right\}  \left\vert \varphi_{j_{1}}%
\cdots\varphi_{j_{P}}\right\rangle \left\langle \varphi_{k_{1}}\cdots
\varphi_{k_{P}}\right\vert \nonumber\\
&  =\binom{N}{P}L_{N}^{P}\left(  \left.  T_{NN}\right\vert _{\mathcal{H}^{N}%
}\right)  =L_{N}^{P}\left(  \Gamma_{N}^{N}\left(  D_{NN}\right)  \right)  \neq
L_{N}^{P}\left(  \Gamma_{P}^{N}\left(  D_{NP}\right)  \right)  .\rfloor
\end{align}


\subsection{The Second Quantization Map $\Gamma.\label{SeCQuant}$}

The restrictions of the field operators give
\begin{align}
\left.  a_{j_{1}}^{\dagger}\cdots a_{j_{N}}^{\dagger}a_{k_{N}}\cdots a_{k_{1}%
}\right\vert _{\mathcal{H}^{N}}  &  =\left\vert \varphi_{j_{1}}\cdots
\varphi_{j_{N}}\right\rangle \left\langle \varphi_{k_{1}}\cdots\varphi_{k_{N}%
}\right\vert \\
\left.  a_{j_{1}}^{\dagger}\cdots a_{j_{P}}^{\dagger}a_{k_{P}}\cdots a_{k_{1}%
}\right\vert _{\mathcal{H}^{N}}  &  =\sum_{1\leq j_{P+1}<\cdots<j_{N}%
\leq\infty}\left\vert \varphi_{j_{1}}\cdots\varphi_{j_{P}}\varphi_{j_{P+1}%
}\cdots\varphi_{j_{N}}\right\rangle \left\langle \varphi_{k_{1}}\cdots
\varphi_{k_{P}}\varphi_{j_{P+1}}\cdots\varphi_{j_{N}}\right\vert ,
\end{align}


The two decompositions of the even Fock algbra%
\begin{equation}
\mathcal{F}_{e}=\mathcal{B}\left(  \mathcal{H}_{\mathcal{F}}\right)  _{e}=%
%TCIMACRO{\dbigoplus \limits_{0\leq N\leq\infty}}%
%BeginExpansion
{\displaystyle\bigoplus\limits_{0\leq N\leq\infty}}
%EndExpansion
\mathcal{B}\left(  \mathcal{H}_{\mathcal{F}}\right)  _{eN}\cong%
%TCIMACRO{\dbigoplus \limits_{0\leq N\leq\infty}}%
%BeginExpansion
{\displaystyle\bigoplus\limits_{0\leq N\leq\infty}}
%EndExpansion
\mathcal{B}\left(  \mathcal{H}^{N}\right)
\end{equation}
can be further refined to give
\begin{equation}
\mathcal{B}\left(  \mathcal{H}^{N}\right)  \cong\mathcal{B}\left(
\mathcal{H}_{\mathcal{F}}\right)  _{eN},
\end{equation}
this isomorphism is a vector space one but not an algebraic one i.e.%
\begin{equation}
\Gamma:\sum_{\substack{1\leq j_{1}<\cdots<j_{P}\leq\infty\\1\leq k_{1}%
<\cdots<k_{P}\leq\infty}}X_{Pj_{1}\cdots j_{P}k_{1}\cdots k_{P}}\left\vert
\varphi_{j_{1}}\cdots\varphi_{j_{P}}\right\rangle \left\langle \varphi_{k_{1}%
}\cdots\varphi_{k_{P}}\right\vert \longrightarrow\sum_{\substack{1\leq
j_{1}<\cdots<j_{P}\leq\infty\\1\leq k_{1}<\cdots<k_{P}\leq\infty}%
}X_{Pj_{1}\cdots j_{P}k_{1}\cdots k_{P}}a_{j_{1}}^{\dagger}\cdots a_{j_{P}%
}^{\dagger}a_{k_{P}}\cdots a_{k_{1}}%
\end{equation}
but
\begin{equation}
\Gamma(XY)\neq\Gamma(X)\Gamma(Y).
\end{equation}


Consider the special case $X_{P}=P_{\mathcal{H}^{P}}$ then
\begin{align}
\Gamma\left(  P_{\mathcal{H}^{P}}\right)   &  =\Gamma\left(  \sum_{1\leq
j_{1}<\cdots<j_{P}\leq\infty}\left\vert \varphi_{j_{1}}\cdots\varphi_{j_{P}%
}\right\rangle \left\langle \varphi_{j_{1}}\cdots\varphi_{j_{P}}\right\vert
\right) \nonumber\\
&  =\sum_{1\leq j_{1}<\cdots<j_{P}\leq\infty}a_{j_{1}}^{\dagger}\cdots
a_{j_{P}}^{\dagger}a_{j_{P}}\cdots a_{j_{1}}=\mathcal{N}_{P},
\end{align}%
\begin{equation}
\frac{\left\langle \Psi_{N}\left.  \mathcal{N}_{P}\right\vert \Psi
_{N}\right\rangle }{\left\langle \left.  \Psi_{N}\right\vert \Psi
_{N}\right\rangle }=\binom{N}{P},
\end{equation}%
\begin{equation}
\Gamma_{P}^{N}\left(  P_{\mathcal{H}^{P}}\right)  =\binom{N}{P}^{-1}\left.
\Gamma\left(  P_{\mathcal{H}^{P}}\right)  \right\vert _{\mathcal{H}^{N}}%
\end{equation}
and
\begin{equation}
P_{\mathcal{H}^{N}}=\Gamma_{P}^{N}\left(  P_{\mathcal{H}^{P}}\right)  .
\end{equation}


As one can expand the second quantized operators in terms of ket-bra's one
obtains
\begin{equation}
\Gamma\left(  \left\vert \varphi_{j}\right\rangle \left\langle \varphi
_{k}\right\vert \right)  =a_{j}^{\dagger}a_{k}=\left\vert \varphi
_{j}\right\rangle \left\langle \varphi_{k}\right\vert +\sum_{1\leq N\leq
\infty}\sum_{1\leq j_{2}<\cdots<j_{N}\leq\infty}\left\vert \varphi_{j}%
\varphi_{j_{2}}\cdots\varphi_{j_{N}}\right\rangle \left\langle \varphi
_{k}\varphi_{j_{2}}\cdots\varphi_{j_{N}}\right\vert .
\end{equation}


Each element of $%
%TCIMACRO{\dbigoplus \limits_{0\leq N\leq\infty}}%
%BeginExpansion
{\displaystyle\bigoplus\limits_{0\leq N\leq\infty}}
%EndExpansion
\mathcal{B}\left(  \mathcal{H}^{N}\right)  $ (which is a direct sum of
algebras) can be represented as $\left(
\begin{array}
[c]{c}%
X_{0}\\
\vdots\\
X_{\infty}%
\end{array}
\right)  _{FQ}$ (with only finitely many non zero elements), $X_{N}%
\in\mathcal{B}\left(  \mathcal{H}^{N}\right)  $. The direct sum $%
%TCIMACRO{\dbigoplus \limits_{0\leq N\leq\infty}}%
%BeginExpansion
{\displaystyle\bigoplus\limits_{0\leq N\leq\infty}}
%EndExpansion
\mathcal{B}\left(  \mathcal{H}_{\mathcal{F}}\right)  _{eN}$ is only a vector
space direct sum and not a direct sum of algebras but we can still represent
each element as $\left(
\begin{array}
[c]{c}%
Y_{0}\\
\vdots\\
Y_{\infty}%
\end{array}
\right)  _{SQ},$ $Y_{N}\in\mathcal{B}\left(  \mathcal{H}_{\mathcal{F}}\right)
_{eN}$. Using the direct sum representations we obtain
\begin{equation}
\Gamma\left(  \left(
\begin{array}
[c]{c}%
X_{0}\\
\vdots\\
X_{\infty}%
\end{array}
\right)  _{FQ}\right)  =\left(
\begin{array}
[c]{c}%
\Gamma\left(  X_{0}\right) \\
\vdots\\
\Gamma\left(  X_{\infty}\right)
\end{array}
\right)  _{SQ}=\left(
\begin{array}
[c]{c}%
Y_{0}\\
\vdots\\
Y_{\infty}%
\end{array}
\right)  _{SQ}%
\end{equation}
and%
\begin{equation}
\Gamma(X_{P})=\sum_{P\leq N\leq\infty}\left.  \Gamma\left(  X_{P}\right)
\right\vert _{\mathcal{H}^{N}}=\sum_{P\leq N\leq\infty}\binom{N}{P}\Gamma
_{P}^{N}\left(  X_{P}\right)  .
\end{equation}
We can reexpress $\Gamma(X_{P})$ wrt the decomposition $%
%TCIMACRO{\dbigoplus \limits_{0\leq N\leq\infty}}%
%BeginExpansion
{\displaystyle\bigoplus\limits_{0\leq N\leq\infty}}
%EndExpansion
\mathcal{B}\left(  \mathcal{H}^{N}\right)  $ as
\begin{equation}
\left(
\begin{array}
[c]{c}%
0\\
\vdots\\
X_{P}\\
\vdots\\
\binom{N}{P}\Gamma_{P}^{N}\left(  X_{P}\right) \\
\vdots
\end{array}
\right)  _{FQ}%
\end{equation}
leading to $\Gamma(\sum\limits_{0\leq P\leq N}X_{P})=\sum\limits_{0\leq P\leq
N}\Gamma(X_{P})$ being given by
\begin{equation}
\Gamma\left(  \left(
\begin{array}
[c]{c}%
X_{0}\\
\vdots\\
X_{N}\\
\vdots\\
X_{\infty}%
\end{array}
\right)  _{FQ}\right)  =\left(
\begin{array}
[c]{c}%
\vdots\\
\sum\limits_{0\leq P\leq N}\binom{N}{P}\Gamma_{P}^{N}\left(  X_{P}\right) \\
\vdots
\end{array}
\right)  _{FQ}.
\end{equation}


\begin{theorem}
$X_{P}\geq0$ $\Rightarrow\left.  \Gamma\left(  X_{P}\right)  \right\vert
_{\mathcal{H}^{N}}\geq0,$ $N>P$

\begin{proof}%
\begin{align}
\left\langle \Psi_{N}\left\vert \left.  \Gamma\left(  X_{P}\right)
\right\vert _{\mathcal{H}^{N}}\right.  \Psi_{N}\right\rangle  &  =\left\langle
\Psi_{N}\left\vert \sum_{\substack{1\leq j_{1}<\cdots<j_{P}\leq\infty\\1\leq
k_{1}<\cdots<k_{P}\leq\infty}}X_{j_{1}\cdots j_{P}k_{1}\cdots k_{P}}a_{j_{1}%
}^{\dagger}\cdots a_{j_{P}}^{\dagger}a_{k_{P}}\cdots a_{k_{1}}\right.
\Psi_{N}\right\rangle \nonumber\\
=  &  \sum_{\substack{1\leq j_{1}<\cdots<j_{P}\leq\infty\\1\leq k_{1}%
<\cdots<k_{P}\leq\infty}}X_{Pj_{1}\cdots j_{P}k_{1}\cdots k_{P}}\left\langle
\Psi_{N}a_{j_{1}}^{\dagger}\cdots a_{j_{P}}^{\dagger}a_{k_{P}}\cdots a_{k_{1}%
}\Psi_{N}\right\rangle \nonumber\\
=\binom{N}{P}  &  \sum_{\substack{1\leq j_{1}<\cdots<j_{P}\leq\infty\\1\leq
k_{1}<\cdots<k_{P}\leq\infty}}X_{Pj_{1}\cdots j_{P}k_{1}\cdots k_{P}}L_{N}%
^{P}\left(  \left\vert \Psi_{N}\right\rangle \left\langle \Psi_{N}\right\vert
\right)  _{k_{1}\cdots k_{P}j_{1}\cdots j_{P}}\nonumber\\
&  \Rightarrow\operatorname{Tr}\left\{  X_{P}L_{N}^{P}\left(  \left\vert
\Psi_{N}\right\rangle \left\langle \Psi_{N}\right\vert \right)  \right\}
\geq0
\end{align}
as $X_{P}\geq0.$
\end{proof}
\end{theorem}

\section{The Kummer Cone}

Cones are defined by
\begin{equation}
C=\left\{  \left.  X\right\vert \alpha X\in X\right\}
\end{equation}
they do not have be convex sets nor linear spaces, if they are they are convex
cones or linear cones.

Let $X_{P}\in\mathcal{B}_{sa}\left(  \mathcal{H}^{P}\right)  $ such
$\sigma\left(  X_{P}\right)  \cap\left(  -\infty,0\right)  \neq\phi$ then
\begin{align}
\inf_{\Phi_{P}\in\mathcal{H}^{P}}\left\{  \frac{\left\langle \left.  \Phi
_{P}\right\vert X_{P}\Phi_{P}\right\rangle }{\left\langle \left.  \Phi
_{P}\right\vert \Phi_{P}\right\rangle }\right\}   &  =\inf_{D^{P}\in
S_{P}\subset\mathcal{B}_{1sa+}\left(  \mathcal{H}^{P}\right)  }\left\{
\operatorname*{Tr}\left\{  X_{P}D^{P}\right\}  \right\} \nonumber\\
&  =\inf_{D^{P}\in S_{P}\subset\mathcal{B}_{1sa+}\left(  \mathcal{H}%
^{P}\right)  }\left\{  \left\langle \left.  X_{P}\right\vert D^{P}%
\right\rangle _{P}\right\}  =\alpha\left(  X_{P}\right)  <0
\end{align}
($\alpha\left(  X_{P}\right)  $ is the lowest eigenvalue or generalized
eigenvalue of $X_{P}$) then by shifting one obtains the positive operator
\begin{equation}
X_{P}-\alpha\left(  X_{P}\right)  P_{P}\geq0\Leftrightarrow X_{P}\geq
\alpha\left(  X_{P}\right)  P_{P},
\end{equation}
which expands to a positive $N$-particle operator%
\begin{equation}
\Gamma_{P}^{N}\left(  X_{P}-\alpha\left(  X_{P}\right)  P_{P}\right)
=\Gamma_{P}^{N}\left(  X_{P}\right)  -\alpha\left(  X_{P}\right)  \Gamma
_{P}^{N}\left(  P_{P}\right)  =\Gamma_{P}^{N}\left(  X_{P}\right)
-\alpha\left(  X_{P}\right)  P_{N}\geq0.
\end{equation}
Consider the expansion of the unshifted operator (which by hypothesis is not
positive) then
\begin{equation}
\inf_{\Phi_{N}\in\mathcal{H}^{N}}\left\{  \frac{\left\langle \left.  \Phi
_{N}\right\vert \Gamma_{P}^{N}\left(  X_{P}\right)  \Phi_{N}\right\rangle
}{\left\langle \left.  \Phi_{N}\right\vert \Phi_{N}\right\rangle }\right\}
=\alpha\left(  \Gamma_{P}^{N}\left(  X_{P}\right)  \right)  \geq\alpha\left(
X_{P}\right)  ,
\end{equation}
as%
\begin{align}
\alpha\left(  \Gamma_{P}^{N}\left(  X_{P}\right)  \right)   &  =\inf_{\Phi
_{N}\in\mathcal{H}^{N}}\left\{  \frac{\left\langle \left.  \Phi_{N}\right\vert
\Gamma_{P}^{N}\left(  X_{P}\right)  \Phi_{N}\right\rangle }{\left\langle
\left.  \Phi_{N}\right\vert \Phi_{N}\right\rangle }\right\}  =\inf_{D^{N}\in
S_{N}}\left\{  \left\langle \left.  \Gamma_{P}^{N}\left(  X_{P}\right)
\right\vert D^{N}\right\rangle _{N}\right\} \nonumber\\
&  =\inf_{D^{N}\in S_{N}}\left\{  \left\langle \left.  X_{P}\right\vert
L_{N}^{P}\left(  D^{N}\right)  \right\rangle _{P}\right\}  =\inf_{D_{NP}%
\in\tilde{S}_{NP}}\left\{  \left\langle \left.  X_{P}\right\vert
D_{NP}\right\rangle _{P}\right\}  \geq\alpha\left(  X_{P}\right)
\end{align}
due to $\tilde{S}_{NP}=L_{N}^{P}\left(  S_{N}\right)  \subset S_{P}$. Thus the
expanded operator is less negative and by shifting
\begin{equation}
\Gamma_{P}^{N}\left(  X_{P}\right)  -\alpha\left(  \Gamma_{P}^{N}\left(
X_{P}\right)  \right)  P_{N}\geq0.
\end{equation}
Hence
\begin{equation}
\alpha\left(  X_{P}\right)  -\alpha\left(  \Gamma_{P}^{N}\left(  X_{P}\right)
\right)  =\inf_{\Phi_{P}\in\mathcal{H}^{P}}\left\{  \frac{\left\langle \left.
\Phi_{P}\right\vert \left\{  X_{P}-\alpha\left(  \Gamma_{P}^{N}\left(
X_{P}\right)  P_{P}\right)  \right\}  \Phi_{P}\right\rangle }{\left\langle
\left.  \Phi_{P}\right\vert \Phi_{P}\right\rangle }\right\}  \leq0
\end{equation}
i.e.%
\begin{equation}
X_{P}-\alpha\left(  \Gamma_{P}^{N}\left(  X_{P}\right)  \right)  P_{P}\ngeq0.
\end{equation}


\begin{theorem}
$\exists Y_{P}=\left.  \Gamma\left(  Y_{P}\right)  \right\vert _{\mathcal{H}%
^{P}}\ngeq0$ $\ni\Gamma_{P}^{N}\left(  Y_{P}\right)  =\binom{N}{P}^{-1}\left.
\Gamma\left(  Y_{P}\right)  \right\vert _{\mathcal{H}^{N}}\geq0,$ $N>P$

\begin{proof}
For $P=1$ let
\begin{equation}
X_{1}=-\left\vert \psi\right\rangle \left\langle \psi\right\vert
\end{equation}
then
\begin{equation}
\inf_{\varphi\in\mathcal{H}^{1}}\left\{  \frac{\left\langle \left.
\varphi\right\vert X_{1}\varphi\right\rangle }{\left\langle \left.
\varphi\right\vert \varphi\right\rangle }\right\}  =\alpha\left(
X_{1}\right)  =-1<0
\end{equation}
and
\begin{align}
&  \inf_{\Phi_{N}\in\mathcal{H}^{N}}\left\{  \frac{\left\langle \left.
\Phi_{N}\right\vert \Gamma_{1}^{N}\left(  X_{1}\right)  \Phi_{N}\right\rangle
}{\left\langle \left.  \Phi_{N}\right\vert \Phi_{N}\right\rangle }\right\}
=\inf_{\Phi_{N}\in\mathcal{H}^{N}}\operatorname{Tr}\left\{  \Gamma_{1}%
^{N}\left(  X_{1}\right)  \frac{\left\vert \Phi_{N}\right\rangle \left\langle
\Phi_{N}\right\vert }{\left\langle \left.  \Phi_{N}\right\vert \Phi
_{N}\right\rangle }\right\} \nonumber\\
&  =\inf_{\Phi_{N}\in\mathcal{H}^{N}}\operatorname{Tr}\left\{  X_{1}L_{N}%
^{1}\left(  \frac{\left\vert \Phi_{N}\right\rangle \left\langle \Phi
_{N}\right\vert }{\left\langle \left.  \Phi_{N}\right\vert \Phi_{N}%
\right\rangle }\right)  \right\}  =\inf_{\Phi_{N}\in\mathcal{H}^{N}%
}\operatorname{Tr}\left\{  \left(  -\left\vert \psi\right\rangle \left\langle
\psi\right\vert \right)  L_{N}^{1}\left(  \frac{\left\vert \Phi_{N}%
\right\rangle \left\langle \Phi_{N}\right\vert }{\left\langle \left.  \Phi
_{N}\right\vert \Phi_{N}\right\rangle }\right)  \right\} \nonumber\\
&  =-\frac{1}{N}=\alpha\left(  \Gamma_{1}^{N}\left(  X_{1}\right)  \right)
>\alpha\left(  X_{1}\right)  .
\end{align}
Thus
\begin{align}
&  \inf_{\varphi\in\mathcal{H}^{1}}\left\{  \frac{\left\langle \left.
\varphi\right\vert \left\{  X_{1}-\alpha\left(  \Gamma_{1}^{N}\left(
X_{1}\right)  \right)  \right\}  \varphi\right\rangle }{\left\langle \left.
\varphi\right\vert \varphi\right\rangle }\right\}  =\alpha\left(
X_{1}\right)  -\alpha\left(  \Gamma_{1}^{N}\left(  X_{1}\right)  \right) \\
&  =-1+\frac{1}{N}=\frac{1-N}{N}<0
\end{align}
and trivally
\begin{align}
&  \inf_{\Phi\in\mathcal{H}^{N}}\left\{  \frac{\left\langle \left.
\Phi\right\vert \Gamma_{1}^{N}\left(  X_{1}-\alpha\left(  \Gamma_{1}%
^{N}\left(  X_{1}\right)  P_{1}\right)  \right)  \Phi\right\rangle
}{\left\langle \left.  \Phi\right\vert \Phi\right\rangle }\right\}  =\\
&  \alpha\left(  \Gamma_{1}^{N}\left(  X_{1}\right)  \right)  -\alpha\left(
\Gamma_{1}^{N}\left(  X_{1}\right)  \right)  =0\\
&  \Leftrightarrow\Gamma_{1}^{N}\left(  X_{1}-\alpha\left(  \Gamma_{1}%
^{N}\left(  X_{1}\right)  \right)  \right)  \geq0.
\end{align}
The above is based on the fact for states $D^{N}\in S_{N}$ one has
\begin{align}
&  \sup_{\varphi_{1}\in\mathcal{H}^{1}}\left\{  \frac{\left\langle \left.
\varphi_{1}\right\vert L_{N}^{1}\left(  D^{N}\right)  \varphi_{1}\right\rangle
}{\left\langle \left.  \varphi_{1}\right\vert \varphi_{1}\right\rangle
}\right\}  =\sup_{\varphi_{1}\in\mathcal{H}^{1}}\left\{  \operatorname{Tr}%
\left\{  \Gamma_{1}^{N}\left(  \frac{\left\vert \varphi_{1}\right\rangle
\left\langle \varphi_{1}\right\vert }{\left\langle \left.  \varphi
_{1}\right\vert \varphi_{1}\right\rangle }\right)  D^{N}\right\}  \right\} \\
&  \left.  =\right.  \sup_{\varphi_{1}\in\mathcal{H}^{1}}\left\{
\operatorname{Tr}\left\{  \binom{N}{1}^{-1}\Gamma\left(  \frac{\left\vert
\varphi_{1}\right\rangle \left\langle \varphi_{1}\right\vert }{\left\langle
\left.  \varphi_{1}\right\vert \varphi_{1}\right\rangle }\right)
D^{N}\right\}  \right\} \\
&  \left.  =\right.  \sup_{\substack{\varphi_{1}\in\mathcal{H}^{1}\\\left\Vert
\varphi_{1}\right\Vert =1}}\left\{  \operatorname{Tr}\left\{  \frac{1}{N}%
\Phi\left(  \varphi_{1}\right)  ^{\dagger}\Phi\left(  \varphi_{1}\right)
D^{N}\right\}  \right\}  =\beta\left(  L_{N}^{1}\left(  D^{N}\right)  \right)
=\frac{1}{N},
\end{align}
which is based on the fact that $\Phi\left(  \varphi_{1}\right)  $ is a
partial isometry ($\Phi\left(  \varphi_{1}\right)  ^{\dagger}\Phi\left(
\varphi_{1}\right)  $ and $\Phi\left(  \varphi_{1}\right)  \Phi\left(
\varphi_{1}\right)  ^{\dagger}$ are idempotent projections) and
\begin{equation}
\left\Vert \Phi\left(  \varphi_{1}\right)  \right\Vert =\left\Vert \Phi\left(
\varphi_{1}\right)  ^{\dagger}\right\Vert =1.
\end{equation}
The corresponding 2-particle operators $\Phi\left(  \varphi_{a}\right)
^{\dagger}\Phi\left(  \varphi_{b}\right)  ^{\dagger}\Phi\left(  \varphi
_{b}\right)  \Phi\left(  \varphi_{a}\right)  =$ $\Phi\left(  \varphi
_{a}\right)  ^{\dagger}\Phi\left(  \varphi_{a}\right)  \Phi\left(  \varphi
_{b}\right)  ^{\dagger}\Phi\left(  \varphi_{b}\right)  $ are also idempotent
and thus projectors hence
\begin{equation}
\left\Vert \Phi\left(  \varphi_{b}\right)  \Phi\left(  \varphi_{a}\right)
\right\Vert =\left\Vert \Phi\left(  \varphi_{a}\right)  ^{\dagger}\Phi\left(
\varphi_{b}\right)  ^{\dagger}\right\Vert =1
\end{equation}
and
\begin{align}
&  \sup_{\varphi_{a}\varphi_{b},\in\mathcal{H}^{1}}\left\{  \frac{\left\langle
\left.  \varphi_{b}\varphi_{a}\right\vert L_{N}^{2}\left(  D^{N}\right)
\varphi_{b}\varphi_{a}\right\rangle }{\left\langle \left.  \varphi_{b}%
\varphi_{a}\right\vert \varphi_{b}\varphi_{a}\right\rangle }\right\} \\
&  \left.  =\right.  \sup_{\varphi_{a}\varphi_{b},\in\mathcal{H}^{1}}\left\{
\operatorname{Tr}\left\{  \binom{N}{2}^{-1}\Gamma\left(  \frac{\left\vert
\varphi_{b}\varphi_{a}\right\rangle \left\langle \varphi_{b}\varphi
_{a}\right\vert }{\left\langle \left.  \varphi_{b}\varphi_{a}\right\vert
\varphi_{b}\varphi_{a}\right\rangle }\right)  D^{N}\right\}  \right\} \\
&  \left.  =\right.  \sup_{\substack{\varphi_{a},\varphi_{b}\in\mathcal{H}%
^{1}\\\left\Vert \varphi_{a}\right\Vert =\left\Vert \varphi_{b}\right\Vert
=1}}\left\{  \operatorname{Tr}\left\{  \frac{2}{N\left(  N+1\right)  }%
\Phi\left(  \varphi_{a}\right)  ^{\dagger}\Phi\left(  \varphi_{b}\right)
^{\dagger}\Phi\left(  \varphi_{b}\right)  \Phi\left(  \varphi_{a}\right)
D^{N}\right\}  \right\}  =\frac{2}{N\left(  N+1\right)  },
\end{align}
however $\left\vert \varphi_{b}\varphi_{a}\right\rangle $ are the set of
decomposable unit vectors in $\mathcal{H}^{2}$ which is a subset of all unit
vectors. This set is given by
\begin{align}
&  \left\vert g\right\rangle =\sum_{1\leq j<k\leq\infty}g_{jk}\left\vert
\varphi_{j}\varphi_{k}\right\rangle \\
&  \sum_{1\leq j<k\leq\infty}\left\vert g_{jk}\right\vert ^{2}=1.
\end{align}
Consider
\begin{align}
&  \sum_{_{\substack{1\leq j_{1}<j_{2}\leq r\\1\leq k_{1}<k_{2}\leq r}%
}}\overline{g}_{j_{1}j_{2}}g_{k_{1}k_{2}}\left\langle \left.  \varphi_{j_{1}%
}\varphi_{j_{2}}\right\vert L_{N}^{2}\left(  D^{N}\right)  \varphi_{k_{1}%
}\varphi_{k_{2}}\right\rangle =\sum_{\substack{1\leq j_{1}<j_{2}\leq r\\1\leq
k_{1}<k_{2}\leq r}}\overline{g}_{j_{1}j_{2}}g_{k_{1}k_{2}}\operatorname{Tr}%
\left\{  a_{j_{1}}^{\dagger}a_{j_{2}}^{\dagger}a_{k_{2}}a_{k_{1}}D^{N}\right\}
\\
&  =\sum_{_{\substack{1\leq j_{1}\leq s\\1\leq k_{1}\leq s}}}c_{j_{1}j_{1}%
+s}c_{k_{1}k_{1}+s}\operatorname{Tr}\left\{  b_{j_{1}}^{\dagger}b_{j_{1}%
+s}^{\dagger}b_{k_{1}+s}b_{k_{1}}D^{N}\right\}
\end{align}
and an extreme geminal
\begin{align}
&  =\frac{1}{s^{2}}\sum_{_{\substack{1\leq j_{1}\leq s\\1\leq k_{1}\leq s}%
}}\operatorname{Tr}\left\{  b_{j_{1}}^{\dagger}b_{j_{1}+s}^{\dagger}%
b_{k_{1}+s}b_{k_{1}}D^{N}\right\} \\
&  =\frac{1}{s^{2}}\sum_{_{\substack{1\leq j_{1}\leq s\\1\leq k_{1}\leq s}%
}}\operatorname{Tr}\left\{  b_{j_{1}}^{\dagger}b_{j_{1}+s}^{\dagger}%
b_{j_{1}+s}b_{j_{1}}D^{N}\right\}  +\frac{2}{s^{2}}\sum_{_{1\leq j_{1}%
<k_{1}\leq s}}\operatorname{Re}\left\{  \operatorname{Tr}\left\{  b_{j_{1}%
}^{\dagger}b_{j_{1}+s}^{\dagger}b_{k_{1}+s}b_{k_{1}}D^{N}\right\}  \right\} \\
&  =\frac{1}{s^{2}}\binom{N}{2}+\frac{2}{s^{2}}\sum_{_{1\leq j_{1}<k_{1}\leq
s}}\operatorname{Re}\left\{  \operatorname{Tr}\left\{  b_{j_{1}}^{\dagger
}b_{j_{1}+s}^{\dagger}b_{k_{1}+s}b_{k_{1}}D^{N}\right\}  \right\}  .
\end{align}
If the extreme geminal was an eigenvector of $L_{N}^{2}\left(  D^{N}\right)  $
then
\begin{align}
\operatorname{Tr}\left\{  b_{j_{1}}^{\dagger}b_{j_{1}+s}^{\dagger}b_{k_{1}%
+s}b_{k_{1}}D^{N}\right\}   &  =0;\text{ }j_{1}\neq k_{1}\\
&  \Rightarrow\sum_{_{1\leq j_{1}<k_{1}\leq s}}\operatorname{Re}\left\{
\operatorname{Tr}\left\{  b_{j_{1}}^{\dagger}b_{j_{1}+s}^{\dagger}b_{k_{1}%
+s}b_{k_{1}}D^{N}\right\}  \right\}  =0
\end{align}
and
\begin{equation}
\left\langle \left.  g\right\vert L_{N}^{2}\left(  D^{N}\right)
g\right\rangle =\frac{1}{s^{2}}\sum_{_{\substack{1\leq j_{1}<j_{2}\leq
r\\1\leq k_{1}<k_{2}\leq r}}}\left\langle \left.  \varphi_{j_{1}}%
\varphi_{j_{2}}\right\vert L_{N}^{2}\left(  D^{N}\right)  \varphi_{k_{1}%
}\varphi_{k_{2}}\right\rangle =\frac{1}{s^{2}}\binom{N}{2}.
\end{equation}

\end{proof}
\end{theorem}

The Kummer cone $\mathcal{K}_{NP}$ is defined as
\begin{equation}
\mathcal{K}_{NP}=\left\{  \left.  X_{P}\right\vert X_{P}\in\mathcal{B}%
_{sa}\left(  \mathcal{H}^{P}\right)  ,\Gamma_{P}^{N}\left(  X_{P}\right)
\geq0\right\}
\end{equation}
it is \emph{crucial to observe} that
\begin{equation}
\mathcal{K}_{NP}\supset\mathcal{B}_{sa+}\left(  \mathcal{H}^{P}\right)  \text{
and }\Gamma_{P}^{N}\left(  \mathcal{K}_{NP}\right)  \subset\mathcal{B}%
_{sa+}\left(  \mathcal{H}^{N}\right)
\end{equation}


\begin{theorem}
(Garrod \& Percus)
\end{theorem}

$D_{P}\in\tilde{S}_{NP}$ $\Longleftrightarrow\left\{  \left\langle \left.
X_{P}\right\vert D_{P}\right\rangle _{P}\geq0\,\,\forall\,\Gamma_{P}%
^{N}\left(  X_{P}\right)  \geq0\right\}  \equiv\left\langle \left.
X_{P}\right\vert D_{P}\right\rangle _{P}\geq0\,\forall X_{P}\in\mathcal{K}%
_{NP}.$

\begin{corollary}
\begin{theorem}
If $D_{P}\notin\tilde{S}_{NP},$ but $D_{P}\in S_{P}$ then $\exists\,X_{P}%
\in\mathcal{K}_{NP}\supset\mathcal{B}_{sa+}\left(  \mathcal{H}^{P}\right)
\ni$
\begin{equation}
\left\langle \left.  X_{P}\right\vert D_{P}\right\rangle _{P}=\left\langle
\left.  X_{P}\right\vert L_{N}^{P}\left(  D_{N}\right)  \right\rangle
_{P}=\left\langle \left.  \Gamma_{P}^{N}\left(  X_{P}\right)  \right\vert
D_{N}\right\rangle _{N}<0.
\end{equation}

\end{theorem}
\end{corollary}

\section{States}

A state is positive normalized linear function acting on $\mathcal{B}\left(
\mathcal{H}_{\mathcal{F}}\right)  $ it is called normal if it can be
identified with an element of the predual $\mathcal{B}_{1}\left(
\mathcal{H}_{\mathcal{F}}\right)  $ and abnormal otherwise i.e. if it belongs
to $\mathcal{B}\left(  \mathcal{H}_{\mathcal{F}}\right)  ^{\ast}%
\ominus\mathcal{B}_{1}\left(  \mathcal{H}_{\mathcal{F}}\right)  .$ If we focus
on expectation values of self adjoint observables one can restrict attention
to
\begin{subequations}
\label{States}%
\begin{align}
f  &  :\mathcal{B}_{sa}\left(  \mathcal{H}_{F}\right)  \rightarrow\mathbb{R}\\
f\left(  X\right)   &  \geq0\forall X\geq0\\
f\left(  I\right)   &  =1\\
f\left(  X\right)   &  =\operatorname{Tr}\left\{  XD\right\}  ;\text{ }D\in
S\subset\mathcal{B}_{1+}\left(  \mathcal{H}_{\mathcal{F}}\right)  \text{ if
}f\text{ is normal,}%
\end{align}
where
\end{subequations}
\begin{equation}
S=\left\{  \left.  D\right\vert D\in\mathcal{B}_{1+}\left(  \mathcal{H}%
_{\mathcal{F}}\right)  ;\operatorname{Tr}\left\{  D\right\}  =1\right\}
\end{equation}
One can consider $\mathcal{B}_{2sa}\left(  \mathcal{H}_{F}\right)  $ as a real
Hilbert space with inner product $\left(  \left.  \bullet\right\vert
\bullet\right)  $ defined as
\begin{equation}
\left(  \left.  X\right\vert Y\right)  =\operatorname{Tr}\left\{  XY\right\}
.
\end{equation}
As any positive element of $\mathcal{B}_{1sa}\left(  \mathcal{H}_{F}\right)  $
can be expressed as the square of elements of $\mathcal{B}_{1sa}\left(
\mathcal{H}_{\mathcal{F}}\right)  $ i.e.%
\begin{equation}
D=Q^{2}%
\end{equation}
the set of states $S$ can be described as
\begin{equation}
S=\left\{  \left.  Q\right\vert \left(  \left.  Q\right\vert Q\right)
=1;Q\in\mathcal{B}_{2sa}\left(  \mathcal{H}_{\mathcal{F}}\right)  \right\}  .
\end{equation}
One should note that the correspondence between $Q$ and $D$ is not $1-1$ as
$D$ has many square roots i.e.
\begin{equation}
Q=\sum\left(  \pm\alpha_{j}\right)  \left\vert \Phi_{j}\right\rangle
\left\langle \Phi_{j}\right\vert ,
\end{equation}
where
\begin{equation}
D=\sum\alpha_{j}^{2}\left\vert \Phi_{j}\right\rangle \left\langle \Phi
_{j}\right\vert .
\end{equation}


\section{Partial States}

A partial state is a postive linear functional defined on a subspace
$\mathcal{W}$ of $\mathcal{B}_{sa}\left(  \mathcal{H}_{F}\right)  $ that
contains the identity $I$ and satisfies eq. $\left(  \ref{States}\right)  $
for elements of the subspace. If the subspace does not contain the identity
then the functional can be extended to a subspace that does if
\begin{equation}
\sup_{\substack{0\leq X\leq I\\X\in\mathcal{W}}}f\left(  X\right)  \leq1
\end{equation}
and%
\begin{equation}
\inf_{\substack{X\geq I\\X\in\mathcal{W}}}f\left(  X\right)  \geq1.
\end{equation}


\subsection{Positive Operators contained within the Unit Ball $\mathfrak{B}%
_{1}\left(  \mathcal{B}\left(  \mathcal{H}_{\mathcal{F}}\right)  \right)  $}

The condition $0$ $\leq X\leq I,$ $X\in\mathcal{B}_{+}\left(  \mathcal{H}%
_{\mathcal{F}}\right)  $ can be expressed as
\begin{equation}
0\leq\left\langle \left.  \Phi\right\vert X\Phi\right\rangle \leq\left\langle
\left.  \Phi\right\vert \Phi\right\rangle \text{ }\forall\Phi\in
\mathcal{H}_{\mathcal{F}}%
\end{equation}
which is equivalent to
\begin{equation}
\sup_{\Phi}\left\{  \frac{\left\langle \left.  \Phi\right\vert X\Phi
\right\rangle }{\left\langle \left.  \Phi\right\vert \Phi\right\rangle
}\right\}  \leq1\Leftrightarrow\left\Vert X\right\Vert \leq1;X\geq0
\end{equation}
and $X\geq I$ is equivalent to
\begin{equation}
\inf_{\Phi}\left\{  \frac{\left\langle \left.  \Phi\right\vert X\Phi
\right\rangle }{\left\langle \left.  \Phi\right\vert \Phi\right\rangle
}\right\}  \geq1\Leftrightarrow\left\Vert X\right\Vert \geq1;X\geq0.
\end{equation}
The extreme points of the unit ball are the unit hypersphere
\begin{equation}
\mathfrak{S}\left(  \mathcal{B}\left(  \mathcal{H}_{\mathcal{F}}\right)
\right)  =\operatorname*{ext}\left\{  \mathfrak{B}_{1}\left(  \mathcal{B}%
\left(  \mathcal{H}_{\mathcal{F}}\right)  \right)  \right\}  =\left\{
X\left\vert \left\Vert X\right\Vert =1,X\in\mathcal{B}\left(  \mathcal{H}%
_{\mathcal{F}}\right)  \right.  \right\}  ,
\end{equation}
and the positive hyperant
\begin{equation}
\mathfrak{S}\left(  \mathcal{B}\left(  \mathcal{H}_{\mathcal{F}}\right)
\right)  _{+}=\mathfrak{S}\left(  \mathcal{B}\left(  \mathcal{H}_{\mathcal{F}%
}\right)  \right)  \cap\mathcal{B}_{+}\left(  \mathcal{H}_{\mathcal{F}%
}\right)  .
\end{equation}
The positive hyperant of the unit ball can be decomposed into the convex sets
$\left\{  \left.  \Lambda_{P}\right\vert 0\leq P\leq\infty\right\}  $ defined by

It is obvious that
\begin{equation}
\lambda\leq I,\lambda\in\Lambda_{P},1\leq P\leq\infty.
\end{equation}
The identity operator $I\in\mathcal{B}\left(  \mathcal{H}_{\mathcal{F}%
}\right)  $ can be expressed as a sum of orthogonal projectors onto the
subspaces $\mathcal{B}\left(  \mathcal{H}^{N}\right)  $ as%
\begin{equation}
I=\sum_{0\leq N\leq\infty}P_{N}%
\end{equation}
thus
\begin{equation}
I>P_{N},\,0\leq N\leq\infty.
\end{equation}
The projectors $\left\{  \left.  P_{N}\right\vert 0\leq N\leq\infty\right\}  $
act as the identity elements $\left\{  \left.  I_{N}\right\vert 0\leq
N\leq\infty\right\}  $ on the spaces $\left\{  \left.  \mathcal{H}%
^{N}\right\vert 0\leq N\leq\infty\right\}  $ i.e.%
\begin{equation}
I_{N}=\left.  I\right\vert _{\mathcal{H}^{N}}=P_{N}.
\end{equation}
One can observe that
\begin{equation}
P_{N}=\left.  \binom{N}{P}^{-1}\mathcal{N}_{P}\right\vert _{\mathcal{H}^{N}%
};0\leq P\leq N
\end{equation}
so in particular
\begin{equation}
I=\sum_{0\leq N\leq\infty}\left.  N^{-1}\mathcal{N}_{1}\right\vert
_{\mathcal{H}^{N}}.
\end{equation}


If $D\in S_{P}$ $\subset\mathcal{B}_{1sa}\left(  \mathcal{H}^{P}\right)  $
then
\begin{align}
0  &  \leq\left\langle \left.  \Phi\right\vert D\Phi\right\rangle
\leq\left\langle \left.  \Phi\right\vert \Phi\right\rangle \text{ }\forall
\Phi\in\mathcal{H}^{P}\\
\operatorname{Tr}\left\{  D\right\}   &  =1\\
\left\langle \left.  \Phi\right\vert D\Phi\right\rangle  &  =0\text{ }%
\forall\Phi\in\mathcal{H}^{M},\text{ }M\neq P
\end{align}
However it only defines a partial state, $D_{NP},$ on $\mathcal{B}_{sa}\left(
\mathcal{H}_{F}\right)  $ if $D_{NP}\in\tilde{S}_{PN}\subset S_{P}$ for $P\leq
N.$

To further analyse these operators one can introduce the following
definitions
\begin{align}
\Lambda_{P}  &  =\left\{  \lambda\left\vert \lambda=\sum_{\substack{1\leq
j_{1}<\cdots<j_{P}\leq\infty\\1\leq k_{1}<\cdots<k_{P}\leq\infty}%
}\lambda_{j_{1}\cdots j_{P}k_{1}\cdots k_{P}}a_{j_{1}}^{\dagger}\cdots
a_{j_{P}}^{\dagger}a_{k_{P}}\cdots a_{k_{1}}\right.  \right. \nonumber\\
&  \left.  \frac{\left\langle \left.  \Phi\right\vert \lambda\Phi\right\rangle
}{\left\langle \left.  \Phi\right\vert \Phi\right\rangle }\leq1\text{ }%
\forall\left\vert \Phi\right\rangle \in\mathcal{H}^{P},\lambda\in
\mathcal{B}_{sa}\left(  \mathcal{H}_{\mathcal{F}}\right)  \right\}
\end{align}
which are convex sets. It is obvious that
\begin{equation}
\lambda\leq I,\lambda\in\Lambda_{P},1\leq P\leq\infty.
\end{equation}


\section{Invariant Subspaces of SU$\left(  \mathcal{H}^{1}\right)  $}

The expansion map is the adjoint of contraction map i.e.
\begin{equation}
\left\langle \left.  \Gamma_{P}^{N}\left(  X\right)  \right\vert
Y\right\rangle _{N}=\left\langle \left.  X\right\vert L_{N}^{P}\left(
Y\right)  \right\rangle _{P}.
\end{equation}
One can expand $\mathcal{B}_{1sa}\left(  \mathcal{H}^{N}\right)  $ as a direct
sum
\begin{equation}
\mathcal{B}_{1sa}\left(  \mathcal{H}^{N}\right)  =%
%TCIMACRO{\dbigoplus \limits_{0\leq P\leq N}}%
%BeginExpansion
{\displaystyle\bigoplus\limits_{0\leq P\leq N}}
%EndExpansion
\mathfrak{B}_{NP},
\end{equation}
where
\begin{equation}
\mathfrak{B}_{NP}=\ker\left\{  L_{N}^{P-1}\right\}  \ominus\ker\left\{
L_{N}^{P}\right\}  ,1\leq P\leq N.
\end{equation}
It is easy to see that%
\begin{equation}
\ker\left\{  L_{N}^{0}\right\}  \supseteq\cdots\supseteq\ker\left\{  L_{N}%
^{N}\right\}
\end{equation}
and%
\begin{equation}
\operatorname*{Tr}\left\{  Y\right\}  =L_{N}^{0}\left(  Y\right)
\end{equation}
so that
\begin{equation}
Y\in\ker\left\{  L_{N}^{0}\right\}  \Leftrightarrow\operatorname*{Tr}\left\{
Y\right\}  =0
\end{equation}
and hence
\begin{equation}
Y\in\ker\left\{  L_{N}^{P}\right\}  ,1\leq P\leq N\Rightarrow
\operatorname*{Tr}\left\{  Y\right\}  =0.
\end{equation}
One should note that
\begin{equation}
\ker\left\{  L_{N}^{N}\right\}  =\phi
\end{equation}
as
\begin{equation}
L_{N}^{N}=\widehat{I_{N}}.
\end{equation}
Any operator $X_{N}\in\mathcal{B}_{1sa}\left(  \mathcal{H}^{N}\right)  $ can
be expanded as
\begin{equation}
X_{N}=X_{NN}+\cdots+X_{N1}+X_{N0},
\end{equation}
where
\begin{equation}
X_{NP}\in\mathfrak{B}_{NP};1\leq P\leq N
\end{equation}
and
\begin{equation}
X_{N0}=X_{N}-X_{NN}+\cdots+X_{N1},
\end{equation}
so that
\begin{align}
\operatorname*{Tr}\left\{  X_{N}\right\}   &  =\operatorname*{Tr}\left\{
X_{N0}\right\} \nonumber\\
\operatorname*{Tr}\left\{  X_{NP}\right\}   &  =0,1\leq P\leq N.
\end{align}


Each subspace $\mathcal{H}^{N}$ carries a representation of $SU\left(
\mathcal{H}^{1}\right)  $ given by
\begin{equation}
u_{N}\left\vert \varphi_{1}\cdots\varphi_{N}\right\rangle =u\wedge\cdots\wedge
u\left\vert \varphi_{1}\cdots\varphi_{N}\right\rangle =\left\vert u\varphi
_{1}\cdots u\varphi_{N}\right\rangle
\end{equation}
and it can be shown that
\begin{equation}
L_{N}^{P}\left(  u_{N}Xu_{N}^{\dagger}\right)  =u_{P}L_{N}^{P}\left(
X\right)  u_{P}^{\dagger},X\in\mathcal{B}_{1sa}\left(  \mathcal{H}^{N}\right)
\end{equation}
thus
\begin{equation}
u_{N}Xu_{N}^{\dagger}\in\ker\left\{  L_{N}^{P}\right\}  \Leftrightarrow
X\in\ker\left\{  L_{N}^{P}\right\}
\end{equation}
i.e. $\ker\left\{  L_{N}^{P}\right\}  $ is invariant (but reducible) under the
action of $SU\left(  \mathcal{H}^{1}\right)  $
\begin{equation}
\widehat{u_{N}}:\ker\left\{  L_{N}^{P}\right\}  \rightarrow\ker\left\{
L_{N}^{P}\right\}
\end{equation}
and
\begin{equation}
\widehat{u_{P}}\circ L_{N}^{P}=L_{N}^{P}\circ\widehat{u_{N}}.
\end{equation}
The Hilbert-Schmidt space $\mathcal{B}_{2sa}\left(  \mathcal{H}^{N}\right)  $
carries a representation of $SU\left(  \mathcal{H}^{1}\right)  $
\begin{equation}
\widehat{u_{N}}\left(  A\right)  =u_{N}Au_{N}^{\dagger},\,A\in\mathcal{B}%
_{2sa}\left(  \mathcal{H}^{N}\right)
\end{equation}
that is reducible and has the direct sum decomposition
\begin{equation}
\mathcal{B}_{2sa}\left(  \mathcal{H}^{N}\right)  =%
%TCIMACRO{\dbigoplus \limits_{0\leq P\leq N}}%
%BeginExpansion
{\displaystyle\bigoplus\limits_{0\leq P\leq N}}
%EndExpansion
\mathfrak{B}_{NP},
\end{equation}
into \emph{orthogonal }irreducible subspaces $\left\{  \mathfrak{B}_{NP};0\leq
P\leq N\right\}  ,$ where
\begin{equation}
\mathfrak{B}_{N0}=\mathbb{R}.
\end{equation}
$.$ These spaces have the property
\begin{align}
X_{N}  &  =\sum_{0\leq P\leq N}X_{NP},X_{NP}\in\mathfrak{B}_{NP};0\leq P\leq
N\\
\sum_{Q+1\leq P\leq N}X_{NP}  &  \in\ker\left\{  L_{N}^{Q}\right\}  ,
\end{align}
where
\begin{equation}
X_{N0}=\operatorname*{Tr}\left\{  X_{N}\right\}  P_{N}.
\end{equation}
If $X_{N}$ $\notin\mathcal{B}_{1sa}\left(  \mathcal{H}^{N}\right)  $ then
$\operatorname*{Tr}\left\{  X_{N}\right\}  $ can be formally set as $\infty$
or $-\infty.$

The pure part expansion can be generalized to the space $\mathcal{B}%
_{2sa}\left(  \mathcal{H}_{F}\right)  _{e}$ which expanded as
\begin{align}
\mathcal{B}_{0sa}\left(  \mathcal{H}_{F}\right)  _{e}  &  =%
%TCIMACRO{\dbigoplus \limits_{0\leq N\leq\infty}}%
%BeginExpansion
{\displaystyle\bigoplus\limits_{0\leq N\leq\infty}}
%EndExpansion
\mathcal{B}_{0sa}\left(  \mathcal{H}^{N}\right)  _{e}=%
%TCIMACRO{\dbigoplus \limits_{\substack{0\leq N\leq\infty\\0\leq P\leq N}}}%
%BeginExpansion
{\displaystyle\bigoplus\limits_{\substack{0\leq N\leq\infty\\0\leq P\leq N}}}
%EndExpansion
\mathfrak{B}_{NP}\\
X  &  =\sum_{\substack{0\leq N\leq\infty\\0\leq P\leq N}}X_{NP}=\sum_{0\leq
P\leq\infty}\sum_{P\leq N\leq\infty}X_{NP}=\sum_{0\leq P\leq\infty}X_{P},
\end{align}
where
\begin{equation}
X_{P}=\sum_{P\leq N\leq\infty}X_{NP}.
\end{equation}
In particular
\begin{equation}
X_{0}=\sum_{0\leq N\leq\infty}X_{N0}=\operatorname*{Tr}\left\{  X\right\}  I
\end{equation}
and
\begin{equation}
X=\sum_{0\leq P\leq\infty}X_{P_{j_{1}\cdots j_{P}k_{1}\cdots k_{P}}}a_{j_{1}%
}^{\dagger}\cdots a_{j_{P}}^{\dagger}a_{k_{P}}\cdots a_{k_{1}}.
\end{equation}
The space $\mathcal{B}_{2sa}\left(  \mathcal{H}_{F}\right)  _{e}$ thus also
carries a representation of $SU\left(  \mathcal{H}^{1}\right)  $as
\begin{align}
\widehat{u}(A)  &  =uAu^{\dagger}=\sum_{0\leq N\leq\infty}\widehat{u_{N}%
}\left(  A_{N}\right)  =\sum_{\substack{0\leq P\leq N\\0\leq N\leq\infty
}}\widehat{u_{N}}\left(  A_{NP}\right)  ,\\
\widehat{u_{N}}  &  =\left.  \widehat{u}\right\vert _{\mathcal{B}_{2sa}\left(
\mathcal{H}^{N}\right)  }%
\end{align}
and
\begin{equation}
u\left(  \lambda\right)  =\exp i\left\{  \sum_{1\leq j,k\leq\infty}%
\lambda_{jk}a_{j}^{\dagger}a_{k}\right\}  ;\lambda_{jk}=\overline{\lambda
}_{kj},1\leq j,k\leq\infty,\operatorname*{Tr}\left\{  \lambda\right\}  =0.
\end{equation}
One can also note that
\begin{equation}
\mathcal{B}_{2sa}\left(  \mathcal{H}_{F}\right)  _{e}=\mathfrak{U}\left(
SU\left(  \mathcal{H}^{1}\right)  \right)  ,
\end{equation}
where $\mathfrak{U}\left(  SU\left(  \mathcal{H}^{1}\right)  \right)  $ is a
representation of the universal covering algebra of $SU\left(  \mathcal{H}%
^{1}\right)  .$

\section{Extra}

Consider
\begin{equation}
T_{NP}=\sum_{\substack{1\leq j_{1}<\cdots<j_{P}\leq\infty\\1\leq k_{1}%
<\cdots<k_{P}\leq\infty}}\operatorname{Tr}\left\{  a_{j_{1}}^{\dagger}\cdots
a_{j_{P}}^{\dagger}a_{k_{P}}\cdots a_{k_{1}}D^{N}\right\}  a_{j_{1}}^{\dagger
}\cdots a_{j_{P}}^{\dagger}a_{k_{P}}\cdots a_{k_{1}}%
\end{equation}
then
\begin{equation}
0\leq\left\langle \left.  \Phi\right\vert T_{NP}\Phi\right\rangle \leq
\binom{N}{P}\left\langle \left.  \Phi\right\vert \Phi\right\rangle \text{
}\forall\Phi\in\mathcal{H}^{P}%
\end{equation}
and for $\Phi\in\mathcal{H}^{P}$
\begin{align}
&  \sum_{\substack{1\leq j_{1}<\cdots<j_{P}\leq\infty\\1\leq k_{1}%
<\cdots<k_{P}\leq\infty}}\operatorname{Tr}\left\{  a_{j_{1}}^{\dagger}\cdots
a_{j_{P}}^{\dagger}a_{k_{P}}\cdots a_{k_{1}}D^{N}\right\}  \left\langle
\left.  \Phi\right\vert a_{j_{1}}^{\dagger}\cdots a_{j_{P}}^{\dagger}a_{k_{P}%
}\cdots a_{k_{1}}\Phi\right\rangle \\
&  =\sum_{\substack{1\leq j_{1}<\cdots<j_{P}\leq\infty\\1\leq k_{1}%
<\cdots<k_{P}\leq\infty}}\operatorname{Tr}\left\{  a_{j_{1}}^{\dagger}\cdots
a_{j_{P}}^{\dagger}a_{k_{P}}\cdots a_{k_{1}}D^{N}\right\}  \bar{C}%
_{j_{1}\cdots j_{P}}C_{k_{1}\cdots k_{P}}\\
&  =\sum_{\substack{1\leq j_{1}<\cdots<j_{P}\leq\infty\\1\leq k_{1}%
<\cdots<k_{P}\leq\infty}}D_{NPj_{1}\cdots j_{P}k_{1}\cdots k_{P}}\bar
{C}_{j_{1}\cdots j_{P}}C_{k_{1}\cdots k_{P}}\geq0.
\end{align}
In particular for $\Phi\in\mathcal{H}^{2}$
\begin{equation}
0\leq\left\langle \left.  \Phi\right\vert T_{N2}\Phi\right\rangle \leq
\binom{N}{2}\left\langle \left.  \Phi\right\vert \Phi\right\rangle \text{
}\forall\Phi\in\mathcal{H}^{2},
\end{equation}%
\begin{align}
\left\langle \left.  \Phi\right\vert T_{N2}\Phi\right\rangle =  &
\sum_{\substack{1\leq j_{1}<j_{2}\leq\infty\\1\leq k_{1}<k_{2}\leq\infty
}}\operatorname{Tr}\left\{  a_{j_{1}}^{\dagger}a_{j_{2}}^{\dagger}a_{k_{2}%
}a_{k_{1}}D^{N}\right\}  \left\langle \left.  \Phi\right\vert a_{j_{1}%
}^{\dagger}a_{j_{2}}^{\dagger}a_{k_{2}}a_{k_{1}}\Phi\right\rangle \\
&  \sum_{\substack{1\leq j_{1}<j_{2}\leq\infty\\1\leq k_{1}<k_{2}\leq\infty
}}\operatorname{Tr}\left\{  a_{j_{1}}^{\dagger}a_{j_{2}}^{\dagger}a_{k_{2}%
}a_{k_{1}}D^{N}\right\}  \left\langle \left.  a_{j_{1}}a_{j_{2}}%
\Phi\right\vert a_{k_{2}}a_{k_{1}}\Phi\right\rangle \\
&  =\sum_{\substack{1\leq j_{1}<j_{2}\leq\infty\\1\leq k_{1}<k_{2}\leq\infty
}}\operatorname{Tr}\left\{  a_{j_{1}}^{\dagger}a_{j_{2}}^{\dagger}a_{k_{2}%
}a_{k_{1}}D^{N}\right\}  \bar{C}_{j_{1}j_{2}}C_{k_{1}k_{2}}\\
&  =\sum_{\substack{1\leq j_{1}<j_{2}\leq\infty\\1\leq k_{1}<k_{2}\leq\infty
}}D_{NPj_{1}j_{2}k_{1}k_{2}}\bar{C}_{j_{1}j_{2}}C_{k_{1}k_{2}}\geq0.
\end{align}
and
\begin{equation}
\left.  T_{N2}\right\vert _{\mathcal{H}^{2}}=\binom{N}{2}L_{N}^{2}\left(
D^{N}\right)
\end{equation}
Let $\Phi\in\mathcal{H}^{P};3\leq P\leq N$ then%
\begin{align}
\left\langle \left.  a_{j_{1}}a_{j_{2}}\Phi\right\vert a_{k_{2}}a_{k_{1}}%
\Phi\right\rangle  &  =\sum_{1\leq m,m^{\prime}\leq\infty}\overline{C}%
_{j_{1}j_{2}m^{\prime}}C_{k_{1}k_{2}m}\left\langle \left.  \varphi_{m^{\prime
}}\right\vert \varphi_{m}\right\rangle =\sum_{1\leq m\leq\infty}\overline
{C}_{j_{1}j_{2}m}C_{k_{1}k_{2}m}=R_{j_{1}j_{2}k_{1}k_{2}}\\
\sum_{\substack{1\leq j_{1}<j_{2}\leq\infty\\1\leq k_{1}<k_{2}\leq\infty
}}\overline{V}_{j_{1}j_{2}}R_{j_{1}j_{2}k_{1}k_{2}}V_{k_{1}k_{2}}  &
=\sum_{1\leq m\leq\infty}\left(  \sum_{1\leq j_{1}<j_{2}\leq\infty}%
\overline{V}_{j_{1}j_{2}}\overline{C}_{j_{1}j_{2}m}\right)  \left(
\sum_{1\leq k_{1}<k_{2}\leq\infty}C_{k_{1}k_{2}m}V_{k_{1}k_{2}}\right)  \geq0
\end{align}%
\begin{align}
\left\langle \left.  \Phi\right\vert T_{N2}\Phi\right\rangle =  &
\sum_{\substack{1\leq j_{1}<j_{2}\leq\infty\\1\leq k_{1}<k_{2}\leq\infty
}}\operatorname{Tr}\left\{  a_{j_{1}}^{\dagger}a_{j_{2}}^{\dagger}a_{k_{2}%
}a_{k_{1}}D^{N}\right\}  \left\langle \left.  \Phi\right\vert a_{j_{1}%
}^{\dagger}a_{j_{2}}^{\dagger}a_{k_{2}}a_{k_{1}}\Phi\right\rangle \\
=  &  \sum_{\substack{1\leq j_{1}<j_{2}\leq\infty\\1\leq k_{1}<k_{2}\leq
\infty}}\operatorname{Tr}\left\{  a_{j_{1}}^{\dagger}a_{j_{2}}^{\dagger
}a_{k_{2}}a_{k_{1}}D^{N}\right\}  \left\langle \left.  a_{j_{1}}a_{j_{2}}%
\Phi\right\vert a_{k_{2}}a_{k_{1}}\Phi\right\rangle \\
&  =\sum_{\substack{1\leq j_{1}<j_{2}\leq\infty\\1\leq k_{1}<k_{2}\leq\infty
}}\operatorname{Tr}\left\{  a_{j_{1}}^{\dagger}a_{j_{2}}^{\dagger}a_{k_{2}%
}a_{k_{1}}\left\vert \Phi\right\rangle \left\langle \Phi\right\vert \right\}
\operatorname{Tr}\left\{  a_{j_{1}}^{\dagger}a_{j_{2}}^{\dagger}a_{k_{2}%
}a_{k_{1}}D^{N}\right\} \\
&  \equiv\sum_{\substack{1\leq j_{1}<j_{2}\leq\infty\\1\leq k_{1}<k_{2}%
\leq\infty}}f_{j_{1}j_{2}k_{1}k_{2}}X_{j_{1}j_{2}k_{1}k_{2}}=f(X);\text{ }%
f\in\mathcal{B}_{1sa}\left(  \mathcal{H}^{2}\right)  ,X\in\mathcal{B}%
_{sa}\left(  \mathcal{H}^{2}\right)  ,
\end{align}
where
\begin{equation}
f_{j_{1}j_{2}k_{1}k_{2}}=\operatorname{Tr}\left\{  a_{j_{1}}^{\dagger}%
a_{j_{2}}^{\dagger}a_{k_{2}}a_{k_{1}}D^{N}\right\}  ;X_{j_{1}j_{2}k_{1}k_{2}%
}=\operatorname{Tr}\left\{  a_{j_{1}}^{\dagger}a_{j_{2}}^{\dagger}a_{k_{2}%
}a_{k_{1}}\left\vert \Phi\right\rangle \left\langle \Phi\right\vert \right\}
.
\end{equation}


\section{Extra}

This can also be seen coordinate wise by considering the special case
$X_{2}=P_{2}$ then
\begin{align}
3\Gamma_{2}^{3}\left(  P_{\mathcal{H}^{2}}\right)   &  =\frac{1}{4\left(
3!\right)  }\sum_{\substack{1\leq j_{1},j_{2}\leq\infty\\1\leq k_{1},k_{2}%
\leq\infty\\1\leq l_{1},l_{2},l_{3}\leq\infty}}\delta_{\left(  j_{1}%
j_{2}\right)  \left(  k_{1}k_{2}\right)  }a_{j_{1}}^{\dagger}a_{j_{2}%
}^{\dagger}a_{k_{2}}a_{k_{1}}\left\vert \varphi_{l_{1}}\varphi_{l_{2}}%
\varphi_{l_{3}}\right\rangle \left\langle \varphi_{l_{1}}\varphi_{l_{2}%
}\varphi_{l_{3}}\right\vert \\
&  =\frac{3}{4\left(  3!\right)  }\sum_{\substack{1\leq j_{1},j_{2}\leq
\infty\\1\leq l_{1},l_{2},l_{3}\leq\infty}}\delta_{\left(  j_{1}j_{2}\right)
\left(  l_{1}l_{2}\right)  }\left\vert \varphi_{j_{1}}\varphi_{j_{2}}%
\varphi_{l_{3}}\right\rangle \left\langle \varphi_{l_{1}}\varphi_{l_{2}%
}\varphi_{l_{3}}\right\vert \\
&  =\frac{\left(  3\right)  \left(  4\right)  }{4\left(  3!\right)  }%
\sum_{1\leq l_{1},l_{2},l_{3}\leq\infty}\left\vert \varphi_{l_{1}}%
\varphi_{l_{2}}\varphi_{l_{3}}\right\rangle \left\langle \varphi_{l_{1}%
}\varphi_{l_{2}}\varphi_{l_{3}}\right\vert \\
&  =\frac{\left(  3\right)  \left(  4\right)  }{4}\sum_{1\leq l_{1}%
<l_{2}<l_{3}\leq\infty}\left\vert \varphi_{l_{1}}\varphi_{l_{2}}\varphi
_{l_{3}}\right\rangle \left\langle \varphi_{l_{1}}\varphi_{l_{2}}%
\varphi_{l_{3}}\right\vert ,
\end{align}
hence%
\begin{equation}
\Gamma_{2}^{3}\left(  P_{\mathcal{H}^{2}}\right)  =P_{\mathcal{H}^{3}}.
\end{equation}
Consider $\left.  \Gamma\left(  X_{P}\right)  \right\vert _{\mathcal{H}^{P+1}%
}=\left(  P+1\right)  \Gamma_{P}^{P+1}\left(  X_{P}\right)  $
\begin{align}
\left.  \Gamma\left(  X_{P}\right)  \right\vert _{\mathcal{H}^{P+1}}  &
=\left(  P+1\right)  \Gamma_{P}^{P+1}\left(  X_{P}\right) \\
&  =\frac{1}{\left(  P!\right)  ^{3}\left(  P+1\right)  }\sum_{\substack{1\leq
j_{1},\cdots,j_{P}\leq\infty\\1\leq k_{1},\cdots,k_{P}\leq\infty\\1\leq
l_{1},\cdots,l_{P+1}\leq\infty}}X_{j_{1}\cdots j_{P}k_{1}\cdots k_{P}}%
a_{j_{1}}^{\dagger}\cdots a_{j_{P}}^{\dagger}a_{k_{P}}\cdots a_{k_{1}%
}\left\vert \varphi_{l_{1}}\cdots\varphi_{l_{P+1}}\right\rangle \left\langle
\varphi_{l_{1}}\cdots\varphi_{l_{P+1}}\right\vert \\
&  =\frac{1}{\left(  P!\right)  ^{2}\left(  P+1\right)  }\sum_{\substack{1\leq
j_{1},\cdots,j_{P}\leq\infty\\1\leq l_{1},\cdots,l_{P+1}\leq\infty}%
}X_{j_{1}\cdots j_{P}l_{1}\cdots l_{P}}\left\vert \varphi_{j_{1}}\cdots
\varphi_{j_{P}}\varphi_{l_{P+1}}\right\rangle \left\langle \varphi_{l_{1}%
}\cdots\varphi_{l_{P}}\varphi_{l_{P+1}}\right\vert \Rightarrow\\
\Gamma_{P}^{P+1}\left(  X_{P}\right)   &  =\sum_{\substack{1\leq j_{1}%
<\cdots,<j_{P}\leq\infty\\1\leq l_{1}<\cdots<l_{P+1}\leq\infty}}X_{j_{1}\cdots
j_{P}l_{1}\cdots l_{P}}\left\vert \varphi_{j_{1}}\cdots\varphi_{j_{P}}%
\varphi_{l_{P+1}}\right\rangle \left\langle \varphi_{l_{1}}\cdots
\varphi_{l_{P}}\varphi_{l_{P+1}}\right\vert .
\end{align}
Consider $\left.  \Gamma\left(  X_{2}\right)  \right\vert _{\mathcal{H}^{3}%
}=3\Gamma_{2}^{3}\left(  X_{2}\right)  $
\begin{align}
\left.  \Gamma\left(  X_{2}\right)  \right\vert _{\mathcal{H}^{3}}  &
=3\Gamma_{2}^{3}\left(  X_{2}\right) \\
&  =\frac{1}{\left(  2\right)  ^{3}3}\sum_{\substack{1\leq j_{1},j_{2}%
\leq\infty\\1\leq k_{1},k_{2}\leq\infty\\1\leq l_{1},l_{2},l_{3}\leq\infty
}}X_{j_{1}j_{2}k_{1}k_{2}}a_{j_{1}}^{\dagger}a_{j_{2}}^{\dagger}a_{k_{2}%
}a_{k_{2}}\left\vert \varphi_{l_{1}}\varphi_{l_{2}}\varphi_{l_{3}%
}\right\rangle \left\langle \varphi_{l_{1}}\varphi_{l_{2}}\varphi_{l_{3}%
}\right\vert \\
&  =\frac{1}{\left(  2\right)  ^{3}3}\sum_{\substack{1\leq j_{1},j_{2}%
\leq\infty\\1\leq l_{1},l_{2},l_{3}\leq\infty}}X_{j_{1}j_{2}l_{1}l_{2}%
}\left\vert \varphi_{j_{1}}\varphi_{j_{2}}\varphi_{l_{3}}\right\rangle
\left\langle \varphi_{l_{1}}\varphi_{l_{2}}\varphi_{l_{3}}\right\vert .
\end{align}
Let
\begin{equation}
X_{P}=\left.  \Gamma\left(  X_{P}\right)  \right\vert _{\mathcal{H}^{P}}%
=\sum_{1\leq j\leq\infty}\alpha_{j}\left\vert \Phi_{Pj}\right\rangle
\left\langle \Phi_{Pj}\right\vert -\sum_{1\leq j\leq\infty}\beta_{j}\left\vert
\Phi_{Pj}\right\rangle \left\langle \Phi_{Pj}\right\vert ;\alpha_{j},\beta
_{j}\geq0\,
\end{equation}
and
\begin{equation}
\left\langle \left.  \Gamma_{P}^{N}\left(  X_{P}\right)  \right\vert
\left\vert \Psi_{N}\right\rangle \left\langle \Psi_{N}\right\vert
\right\rangle _{N}=\left\langle \left.  X_{P}\right\vert L_{N}^{P}\left(
\left\vert \Psi_{N}\right\rangle \left\langle \Psi_{N}\right\vert \right)
\right\rangle _{P}\geq0\,\forall\Psi_{N}.
\end{equation}
Then
\begin{align}
&  \left\langle \left.  \left\{  \sum_{1\leq j\leq\infty}\alpha_{j}\left\vert
\Phi_{Pj}^{+}\right\rangle \left\langle \Phi_{Pj}^{+}\right\vert -\sum_{1\leq
j\leq\infty}\beta_{j}\left\vert \Phi_{Pj}^{-}\right\rangle \left\langle
\Phi_{Pj}^{-}\right\vert \right\}  \right\vert L_{N}^{P}\left(  \left\vert
\Psi_{N}\right\rangle \left\langle \Psi_{N}\right\vert \right)  \right\rangle
_{P}=\\
&  \left\langle \left.  \Gamma_{P}^{N}\left(  \sum_{1\leq j\leq\infty}%
\alpha_{j}\left\vert \Phi_{Pj}^{+}\right\rangle \left\langle \Phi_{Pj}%
^{+}\right\vert \right)  \right\vert \left\vert \Psi_{N}\right\rangle
\left\langle \Psi_{N}\right\vert \right\rangle _{N}-\left\langle \left.
\Gamma_{P}^{N}\left(  \sum_{1\leq j\leq\infty}\beta_{j}\left\vert \Phi
_{Pj}^{-}\right\rangle \left\langle \Phi_{Pj}^{-}\right\vert \right)
\right\vert \left\vert \Psi_{N}\right\rangle \left\langle \Psi_{N}\right\vert
\right\rangle _{N}\\
&  \qquad\qquad\qquad\qquad\qquad\qquad\qquad\qquad\qquad\left.  \geq\right.
0\qquad\forall\Psi_{N},
\end{align}
or more generally
\begin{equation}
\left\langle \left.  \Gamma_{P}^{N}\left(  \sum_{1\leq j\leq\infty}\alpha
_{j}\left\vert \Phi_{Pj}^{+}\right\rangle \left\langle \Phi_{Pj}%
^{+}\right\vert \right)  \right\vert \left\vert \Psi_{N}\right\rangle
\left\langle \Psi_{N}\right\vert \right\rangle _{N}\geq\left\langle \left.
\Gamma_{P}^{N}\left(  \sum_{1\leq j\leq\infty}\beta_{j}\left\vert \Phi
_{Pj}^{-}\right\rangle \left\langle \Phi_{Pj}^{-}\right\vert \right)
\right\vert \left\vert \Psi_{N}\right\rangle \left\langle \Psi_{N}\right\vert
\right\rangle _{N}\forall\Psi_{N}%
\end{equation}


\begin{theorem}
\begin{proof}
Is
\begin{equation}
\Gamma_{P}^{N}\left(  \left\vert \Phi_{Pj}^{+}\right\rangle \left\langle
\Phi_{Pj}^{+}\right\vert \right)  \bot\Gamma_{P}^{N}\left(  \left\vert
\Phi_{Pj}^{-}\right\rangle \left\langle \Phi_{Pj}^{-}\right\vert \right)
\end{equation}
i.e.%
\begin{align}
\left(  \Gamma_{P}^{N}\left(  \left\vert \Phi_{Pj}^{+}\right\rangle
\left\langle \Phi_{Pj}^{+}\right\vert \right)  \left\vert \Gamma_{P}%
^{N}\left(  \left\vert \Phi_{Pj}^{-}\right\rangle \left\langle \Phi_{Pj}%
^{-}\right\vert \right)  \right.  \right)  _{N}  &  =\left(  L_{N}^{P}%
\Gamma_{P}^{N}\left(  \left\vert \Phi_{Pj}^{+}\right\rangle \left\langle
\Phi_{Pj}^{+}\right\vert \right)  \left\vert \left\vert \Phi_{Pj}%
^{-}\right\rangle \left\langle \Phi_{Pj}^{-}\right\vert \right.  \right)
_{P}\\
&  =\left(  \left\vert \Phi_{Pj}^{+}\right\rangle \left\langle \Phi_{Pj}%
^{+}\right\vert \left\vert L_{N}^{P}\Gamma_{P}^{N}\left(  \left\vert \Phi
_{Pj}^{-}\right\rangle \left\langle \Phi_{Pj}^{-}\right\vert \right)  \right.
\right)  _{P}=0
\end{align}
$\lceil$yes, think of $AA^{\dagger}$ and $A^{\dagger}A$ if $A$ is compact$,$%
\begin{align}
A  &  =\sum_{1\leq j\leq\infty}\alpha_{j}\left\vert f_{j}\right\rangle
\left\langle e_{j}^{d}\right\vert ;\,A^{\dagger}=\sum_{1\leq j\leq\infty
}\overline{\alpha}_{j}\left\vert e_{j}\right\rangle \left\langle f_{j}%
^{d}\right\vert \\
A  &  :E\rightarrow F;\quad A^{\dagger}:F^{d}\rightarrow E^{d}\\
\text{Let }F  &  \subset F^{d}\text{ and }E\subset E^{d}\\
AA^{\dagger}  &  =\sum_{1\leq j\leq\infty}\left\vert \alpha_{j}\right\vert
^{2}\left\vert f_{j}\right\rangle \left\langle f_{j}^{d}\right\vert
\,;\,A^{\dagger}A=\sum_{1\leq j\leq\infty}\left\vert \alpha_{j}\right\vert
^{2}\left\vert e_{j}\right\rangle \left\langle e_{j}^{d}\right\vert \\
AA^{\dagger}  &  :F\rightarrow F;\quad A^{\dagger}A:E\rightarrow E,
\end{align}
where $\left\{  \left\vert e_{j}\right\rangle ,\left\langle e_{j}%
^{d}\right\vert \right\}  $ and $\left\{  \left\vert f_{j}\right\rangle
,\left\langle f_{j}^{d}\right\vert \right\}  $ are biorthogonal bases. So
\begin{align}
\left\langle \left.  e_{k}\right\vert e_{l}\right\rangle _{E}  &  =\delta
_{kl}\\
\left\langle \left.  Ae_{k}\right\vert Ae_{l}\right\rangle _{F}  &
=\left\langle \left.  e_{k}\right\vert A^{\dagger}Ae_{l}\right\rangle
_{E}=\left\langle e_{k}\left\vert \sum_{1\leq j\leq\infty}\left\vert
\alpha_{j}\right\vert ^{2}\left\vert e_{j}\right\rangle \left\langle e_{j}%
^{d}\right\vert e_{l}\right.  \right\rangle _{E}\\
&  =\delta_{kl}.\rfloor
\end{align}%
\begin{equation}
a_{j}^{\dagger}a_{k}=\left(  \sum_{0\leq N\leq\infty}P_{N}\right)
a_{j}^{\dagger}a_{k}\left(  \sum_{0\leq N\leq\infty}P_{N}\right)  =\sum_{0\leq
N\leq\infty}P_{N}a_{j}^{\dagger}a_{k}P_{N},
\end{equation}
where%
\begin{equation}
\sum_{0\leq N\leq\infty}P_{N}=I_{F}.
\end{equation}

\end{proof}
\end{theorem}


\end{document}